% ============================================================
% T-DOME Paper III: The Loop
% T-DOME three-paper series, Paper III -- October 2026 revision
% ============================================================

\documentclass[12pt]{article}
\usepackage[margin=1in]{geometry}
\usepackage{amsmath,amssymb,amsthm}
\usepackage[hidelinks]{hyperref}
\hypersetup{
  pdftitle={Keeping a Self-Model Calibrated: Drift Detection, the Price of Keeping Up, and Reliability in a Closed Loop},
  pdfauthor={Sidong Liu}
}
\usepackage{booktabs}
\usepackage{physics}
\usepackage{array}
\newcolumntype{P}[1]{>{\raggedright\arraybackslash}p{#1}}

\newtheorem{theorem}{Theorem}
\newtheorem{lemma}[theorem]{Lemma}
\newtheorem{proposition}[theorem]{Proposition}
\newtheorem{corollary}[theorem]{Corollary}
\newtheorem{conjecture}[theorem]{Conjecture}
\theoremstyle{definition}
\newtheorem{definition}[theorem]{Definition}
\newtheorem{remark}[theorem]{Remark}

\newcommand{\E}{\mathbb E}

\title{Keeping a Self-Model Calibrated:\\
Drift Detection, the Price of Keeping Up,\\
and Reliability in a Closed Loop}

\author{Sidong Liu, PhD \\
\small iBioStratix Ltd \\
\small \texttt{sidongliu@hotmail.com}}

\date{Revised October 2026}

\begin{document}
\emergencystretch=2em
\raggedbottom
\maketitle

% ============================================================
\begin{abstract}
A learned instrument model helps an observer decide which parts of its own access are worth maintaining. We study this decision together with drift detection and paid calibration. Smooth even record laws delay linear-drift detection to order $\varepsilon^{-4/5}$, compared with $\varepsilon^{-2/3}$ for first-order probes. For scalar Gaussian transitions, no free background information and quadratic loss, the all-policy long-run cost lies between $C-\tfrac12k_BTG\nu^2$ and $C$, with $C\propto\nu^{2/3}$. The bounds have the same slow-drift leading term; the upper bound uses a Gaussian prior and reads with controllable continuous precision, while integer reads inherit the relaxation lower bound. In an acting body, stationary correlations identify gain and actuator success, and labelled standards supply orientation. Adaptive reference timing has a strict advantage and an exact optimum in a finite-rate example. When the body's working state gates reading and processing, latched repair changes future access and the kernel between acquired records. Retained body history raises task value and, under positive gain and sign drift, can reverse optimal repair at equal sign confidence. Two posterior moments determine terminal access value; continuing enquiry can require a third moment and the joint belief. In a two-observer quadratic task, estimation-optimal and task-optimal authority weights differ. These constructions use one acquired model for task permission, future enquiry and repair under declared physical supplies.
\end{abstract}

\medskip
\noindent\textbf{Keywords}: self-calibration, self-models, Fisher information, change detection, posterior Cram\'er--Rao bound, drift tracking, information thermodynamics, feedback control

% ============================================================
\section{Introduction}
\label{sec:intro}

\subsection{The question}
\label{subsec:question}

Paper~I bounds the usable value of a stored record by $k_BT\,I(X{:}M)$ and associates every feedback protocol with an implicit matching distribution~\cite{Liu2026TDOME_I}. This distribution need not be represented inside the agent. When an agent does represent its own readout channel and decodes with a wrong model, it loses part of the record's value and can fall below an agent without memory. Paper~II studies a retained self-frame that selects reads under a channel budget; in-frame records cannot locate its missing task direction~\cite{Liu2026TDOME_II}. Here such a represented model is checked against an external reference or the pointer value selected by the body's working state. The later body constructions show how acquired records form that model and change its next operations.

We ask how quickly drift is detected, what calibration costs, and how the retained model changes action when the body's function determines future access.

\subsection{Position within the series}
\label{subsec:sequence}

This paper completes the T-DOME (Thermodynamic Dynamics of Observer-Memory Entanglement) series, the third line of a programme whose other lines are HAFF and Q-RAIF.



The three T-DOME papers follow a route that the programme calls memory--self--\allowbreak introspection. An agent keeps records of its world (Paper~I), organizes its access to them through a model of itself (Paper~II), and checks that model against references while it acts (this paper). The links share their objects. The gain--bias channel whose calibrated loss Paper~I computes is the instrument of Section~\ref{sec:loop_world}; the even residual law that Paper~II derives from a budgeted frame is the detection model of Section~\ref{sec:detection}; the conditional drift law cited in Paper~I is proved in Section~\ref{sec:price}. Each link is established in a stated model.

\subsection{Main results}
\label{subsec:results}

The analytical chain runs from detection (Section~\ref{sec:detection})
through all-policy Gaussian maintenance costs (Section~\ref{sec:price})
to tracking (Section~\ref{sec:loop}) and implementation
(Appendix~\ref{sec:cost}).
The joint body model of Section~\ref{sec:autonomous} supplies the main
acting example. Section~\ref{subsec:body_access} then makes working
access conditional on the observer's body: actual history changes
harvest, permission to act, and paid repair. Sign confidence is not
a sufficient repair statistic. Terminal access has a two-moment
value, while continued enquiry generally needs the joint belief.
The older numerical checks and gain--bias bank models are retained
in Appendices~\ref{sec:example}--\ref{sec:loop_world}; the finite-filter
implementation certificates are in Appendix~\ref{app:filter_certificates}.

\subsection{Relation to internal selection}
\label{subsec:haff_g}

Paper~G of HAFF gives a typed model of internal selection: a program register holds a choice, a fixed executor turns program states into couplings or readouts, records are outcomes of instruments, and a fixed rule rewrites the program from the records~\cite{Liu2026HAFF_G}. A finite system can select its own readouts and revise that selection. Within this target/program architecture, a selected target operation is not a complete implementation of the larger executor acting on both registers; a stored description or an additional selector is not excluded. The calibration loop uses this architecture. The self-model is the program, decoding and command rules are the executor, and calibration is the revision rule. History changes the program and the operations it selects. The executor and the labelled standard specify how those changes are used.

\subsection{Related work}
\label{subsec:related}

The Fisher information metric was introduced by Rao~\cite{Rao1945} and characterized by \v{C}encov~\cite{Cencov1982}; Amari developed the natural gradient and information geometry~\cite{Amari1998,AmariNagaoka2000}. The Bayesian Cram\'er--Rao (van Trees) inequality~\cite{vanTrees1968} and its recursive form for filtering, the posterior Cram\'er--Rao bound of Tichavsk\'y, Muravchik and Nehorai~\cite{Tichavsky1998}, are the estimation tools of Sections~\ref{sec:SRCR}--\ref{sec:price}; the Kalman filter~\cite{Kalman1960} attains the latter in the linear Gaussian case. The trade-off between noise and lag in adaptive estimation, and the persistent-excitation condition, are standard in adaptive control~\cite{AstromWittenmark1995}. Thermodynamic length and minimum-dissipation protocols are due to Crooks and Sivak~\cite{Crooks2007,SivakCrooks2012}, and Fisher information enters entropy production through thermodynamic uncertainty relations and information geometry~\cite{BaratoSeifert2015,Ito2018}. Second-order cybernetics, the control of the observing system itself, goes back to von Foerster, building on Ashby's cybernetics~\cite{Ashby1956,vonFoerster2003}. The data-processing bound for hypothesis tests used in Section~\ref{sec:detection} is standard~\cite{CoverThomas2006}. We do not claim that the telescoping inequality behind Theorem~\ref{thm:driftlaw} is new to filtering theory; its use here is to settle the drift law of Paper~I for all policies and to give it a thermodynamic reading.

% ============================================================
\section{Setting}
\label{sec:prelim}

\subsection{Inputs from Papers I and II}
\label{subsec:from_I}\label{subsec:from_II}
Paper~I supplies the record's matched work value, the posterior mismatch
loss and its calibrated expansion~\cite{Liu2026TDOME_I}.
Its gain--bias example has leading loss
$k_BT(1+g^2/(1-b^2))/(4n)$ after $n$ add-half reference reads per sign.
Paper~II supplies a retained foreground, an in-frame blind spot and,
for its two-channel reader, the even information law~\cite{Liu2026TDOME_II}
\begin{equation}\label{eq:fisher_even_II}
 F(\theta)=\frac{\rho^2\sin^2\theta\,(\rho^2+s^2+\rho^2\cos^2\theta)}
 {(s^2+\rho^2\sin^2\theta)^2}.
\end{equation}
An off-frame read resolves its sign fibre. The Gaussian drift law cited
by Paper~I is proved here; its assumptions do not identify either reader
with the bounded acting body below.

\subsection{Fisher information and the Cram\'er--Rao bounds}
\label{subsec:fisher}

\begin{definition}[Fisher information]
\label{def:fisher}
Let $\{p(x\mid\theta):\theta\in\Theta\subset\mathbb R^d\}$ be a family of densities for which differentiation and integration may be interchanged. Its Fisher information matrix is
\begin{equation}
\label{eq:fisher_def}
g_{ij}(\theta)
:=\E_\theta\bigl[\partial_i\log p(x\mid\theta)\,\partial_j\log p(x\mid\theta)\bigr]
=-\E_\theta\bigl[\partial_i\partial_j\log p(x\mid\theta)\bigr].
\end{equation}
Where it is positive definite it is a Riemannian metric on $\Theta$.
\end{definition}

By \v{C}encov's theorem, on finite probability simplices the Fisher metric is, up to a constant factor, the unique Riemannian metric invariant under statistically sufficient embeddings~\cite{Cencov1982}. This makes it the canonical local measure of distinguishability; it does not make a Fisher-based detector physically mandatory, and it does not choose the agent's likelihood or objective. The metric is the Hessian of relative entropy,
\begin{equation}
\label{eq:fisher_KL}
D\bigl(p_\theta\,\Vert\,p_{\theta+d\theta}\bigr)
=\tfrac12\,g_{ij}(\theta)\,d\theta^i\,d\theta^j+O(|d\theta|^3),
\end{equation}
and it bounds estimation: for an unbiased estimator from $n$ independent observations, $\operatorname{Cov}(\hat\theta)\succeq g(\theta)^{-1}/n$~\cite{Rao1945,CoverThomas2006}. For a correlated stream use the Fisher information of its joint likelihood. A decorrelation time alone does not justify replacing the sample count by $T/\tau_E$.

\subsection{Standing assumptions}
\label{subsec:standing_C}

\begin{definition}[Standing assumptions]
\label{def:assumptions_C}
Sections~\ref{sec:detection}--\ref{sec:loop} assume the following. Section~\ref{sec:loop_world} states its own inputs.
\begin{enumerate}
\item[\textup{(C1)}] \textbf{Self-model.} The agent's model of itself is a parameter $\vartheta$ in a finite-dimensional manifold $\Theta$: a frame of Paper~II, the gain and bias of an instrument, or a scalar mismatch coordinate. The family of records $p(e\mid\vartheta)$ that the agent uses to assess it is specified.
\item[\textup{(C2)}] \textbf{Drift.} The parameter that best describes the instrument, $\vartheta^*(t)$, moves. Each result states its drift law: linear, spherical Brownian, a Gaussian random walk, or bounded speed.
\item[\textup{(C3)}] \textbf{Regularity.} Fisher informations are finite and differentiation under the integral sign is allowed to the order used. Uniform lower bounds on information are assumed separately where they are needed.
\end{enumerate}
\end{definition}

% ============================================================
\section{Detecting drift from one's own records}
\label{sec:detection}

A self-model goes stale when the instrument it describes changes. The agent's first question is whether it can notice this from the records it already gathers. Paper~II shows that records read through a frame depend on the frame's misalignment only through an even function of it. This section works out what evenness implies.

\subsection{Even models}
\label{subsec:first_order}

\begin{proposition}[Sign non-identifiability]
\label{prop:first_order_nonidentifiability}
If a record family satisfies $p(e\mid\theta)=p(e\mid-\theta)$ for all $\theta$, then the joint law of any number of conditionally independent records is invariant under $\theta\mapsto-\theta$, and no statistic of the records distinguishes $\theta$ from $-\theta$.
\end{proposition}

\begin{proof}
The joint likelihood is a product of factors that are each invariant under the sign change.
\end{proof}

Paper~II derives this structure instead of assuming it: a frame that reads one of two channels sees an in-frame law that depends on the angle $\theta$ to the coupling only through $\cos\theta$, and the sign is the unobservable fibre of its blind spot~\cite{Liu2026TDOME_II}. Evenness has a second consequence, which governs how fast drift can be noticed at all.

\begin{proposition}[Information and relative entropy of an even model]
\label{prop:even}
Let $p(e\mid\theta)$ be even in $\theta$ and smooth enough near $\theta=0$ that the expansions below may be integrated term by term (this holds for the Gaussian families of Corollary~\ref{cor:three_records}). Write $\ell(\theta)=\log p(e\mid\theta)$ and
\[
\kappa:=\E_0\bigl[\ell''(0)^2\bigr].
\]
Assume $\kappa>0$ for the sharp detection exponents below. Then the score vanishes at $\theta=0$, so the Fisher information $\mathcal I_1(0)$ of one record is zero, and
\begin{equation}
\label{eq:even_expansions}
\mathcal I_1(\theta)=\kappa\,\theta^2+O(\theta^4),\qquad
D\bigl(p_\theta\,\Vert\,p_0\bigr)=\frac{\kappa}{8}\,\theta^4+O(\theta^6).
\end{equation}
\end{proposition}

\begin{proof}
For each $e$ the function $\ell$ is even, so $\ell'(0)=\ell'''(0)=0$ and $\ell'(\theta)=\theta\,\ell''(0)+O(\theta^3)$. Since $p_\theta-p_0=O(\theta^2)$, $\mathcal I_1(\theta)=\E_\theta[\ell'(\theta)^2]=\theta^2\E_0[\ell''(0)^2]+O(\theta^4)$. For the relative entropy write $p_\theta/p_0=1+\eta$ with $\eta=\tfrac12\theta^2H+O(\theta^4)$ and $H=p_0''/p_0=\ell''(0)+\ell'(0)^2=\ell''(0)$. Normalization gives $\E_0\eta=0$, and $(1+\eta)\ln(1+\eta)=\eta+\tfrac12\eta^2+O(\eta^3)$, so $D(p_\theta\Vert p_0)=\tfrac12\E_0\eta^2+O(\theta^6)=\tfrac18\theta^4\E_0H^2+O(\theta^6)$.
\end{proof}

\begin{corollary}[Three record families]
\label{cor:three_records}
\begin{enumerate}
\item[(a)] \emph{Gaussian mean model}, $e\mid\theta\sim\mathcal N(a\theta^2,\sigma_e^2)$: $\mathcal I_1(\theta)=4a^2\theta^2/\sigma_e^2$ and $D(p_\theta\Vert p_0)=a^2\theta^4/(2\sigma_e^2)$ exactly, so $\kappa=4a^2/\sigma_e^2$.
\item[(b)] \emph{Paper~II's in-frame record} $(v\cdot x,y)$ with $D=2$: $\mathcal I_1=F(\theta)$ of~\eqref{eq:fisher_even_II}, and
\begin{equation}
\label{eq:kl_inframe}
D(p_\theta\Vert p_0)=\frac{\rho^2(1-\cos\theta)}{s^2}-\frac12\ln\Bigl(1+\frac{\rho^2\sin^2\theta}{s^2}\Bigr),\qquad
\kappa=\frac{\rho^2(2\rho^2+s^2)}{s^4}.
\end{equation}
\item[(c)] \emph{An off-frame probe record} $(v_\perp\cdot x,y)$: $D(p_\theta\Vert p_0)=-\tfrac12\ln\bigl(1-\rho^2\sin^2\theta/(\rho^2+s^2)\bigr)$, and its Fisher information
\begin{equation}
\label{eq:fisher_probe}
F_\perp(\theta)=\frac{\rho^2\cos^2\theta\,(\rho^2\sin^2\theta+\rho^2+s^2)}{(\rho^2+s^2-\rho^2\sin^2\theta)^2}
\end{equation}
is positive at calibration, $F_\perp(0)=\rho^2/(\rho^2+s^2)$.
\end{enumerate}
\end{corollary}

\begin{proof}
(a) is the normal location family with mean $a\theta^2$. For (b) and (c) the record is a zero-mean Gaussian pair with covariance $\Sigma_\theta=\bigl(\begin{smallmatrix}1&c(\theta)\\c(\theta)&V\end{smallmatrix}\bigr)$, $V=\rho^2+s^2$, where $c=\rho\cos\theta$ in (b) and $c=\rho\sin\theta$ in (c). Then $D(p_\theta\Vert p_0)=\tfrac12[\tr(\Sigma_0^{-1}\Sigma_\theta)-2-\ln(\det\Sigma_\theta/\det\Sigma_0)]$ gives the two relative entropies, and the Fisher information of a Gaussian family with parameter-dependent covariance, $\tfrac12\tr(\Sigma^{-1}\Sigma'\Sigma^{-1}\Sigma')=c'^2(c^2+V)/(V-c^2)^2$, gives~\eqref{eq:fisher_even_II} and~\eqref{eq:fisher_probe}. Expanding~\eqref{eq:kl_inframe} to fourth order gives $\kappa/8=\rho^2(2\rho^2+s^2)/(8s^4)$, in agreement with $F(\theta)/\theta^2\to\rho^2(2\rho^2+s^2)/s^4$.
\end{proof}

Model (a) is a convenient assumption; model (b) is what a budgeted frame actually produces. Proposition~\ref{prop:even} shows that the quadratic onset of information does not depend on the choice: evenness gives the expansion; a nonzero $\kappa$ gives its quadratic onset. The constant $\kappa$ is the only feature of the model that enters the leading order. For records at rate $r$ under linear drift $\theta(t)=\theta_0+\varepsilon t$, the accumulated local information $J_T:=\sum_{t_m\le T}\mathcal I_1(\theta(t_m))\simeq r\int_0^T\mathcal I_1(\theta_0+\varepsilon t)\,dt$ is
\begin{equation}
\label{eq:accumulated_fisher}
J_T=\frac{\kappa r\varepsilon^2T^3}{3}+O(\varepsilon^4T^5)\quad(\theta_0=0),\qquad
J_T=\frac{4a^2r}{3\sigma_e^2\varepsilon}\bigl[(\theta_0+\varepsilon T)^3-\theta_0^3\bigr]\quad\text{in model (a)}.
\end{equation}
It enters the estimation bound of Section~\ref{sec:SRCR}. Detection is governed by a different quantity, the accumulated relative entropy.

\subsection{How long detection takes}
\label{subsec:detection_time}

Fix a horizon $T$ and let records arrive at rate $r$, conditionally independent given the mismatch path $\theta_m=\theta(m/r)$. A test is a function of the records up to time $T$ with the values ``calibrated'' and ``drifted''. It has errors $(\alpha,\beta)$ if it says ``drifted'' with probability at most $\alpha$ when $\theta_m=0$ for all $m$, and ``calibrated'' with probability at most $\beta$ under the drifted law; for a random drift path the drifted law is the mixture over paths. Write $d(p\Vert q)=p\ln\frac pq+(1-p)\ln\frac{1-p}{1-q}$.

\begin{proposition}[Detection time]
\label{prop:detect}
Every test with errors $(\alpha,\beta)$, $\alpha+\beta<1$, satisfies
\begin{equation}
\label{eq:detect_bound}
\sum_{m\le rT}\E\,D\bigl(p_{\theta_m}\,\Vert\,p_0\bigr)\ \ge\ d(1-\beta\,\Vert\,\alpha),
\end{equation}
with the expectation over the drift path if it is random. Consequently, as the drift rate tends to zero:
\begin{enumerate}
\item[(a)] for an even model with $D(p_\theta\Vert p_0)\le\kappa\theta^4/8+C\theta^6$ for all $\theta$ and linear drift $\theta(t)=\varepsilon t$, $T\ge\bigl(40\,d/(\kappa r\varepsilon^4)\bigr)^{1/5}(1+o(1))$;
\item[(b)] if probe positions are fixed in advance with counting function $N_\perp(t)=prt+o(t)$ for $p>0$, their record law has $D(p_\theta\Vert p_0)\le\tfrac12F_1\theta^2+C\theta^4$ with $F_1>0$ its Fisher information at calibration, as off-frame probes do, and the others are as in (a), $T\ge\bigl(6\,d/(pF_1r\varepsilon^2)\bigr)^{1/3}(1+o(1))$;
\item[(c)] for Paper~II's reader with $D=2$ under spherical Brownian drift with generator $\sigma\Delta$, $T\ge\bigl(2\,d/(\kappa r\sigma^2)\bigr)^{1/3}(1+o(1))$,
\end{enumerate}
where $d=d(1-\beta\Vert\alpha)$.
\end{proposition}

\begin{proof}
Let $P_0$ and $P_1$ be the laws of the records up to $T$ under the two hypotheses. The test's decision is a function of the records, so by the data-processing inequality $D(P_1\Vert P_0)$ is at least the relative entropy between the two laws of the decision, which are Bernoulli with parameters $p_1\ge1-\beta$ and $p_0\le\alpha$; since $d(p_1\Vert p_0)$ increases in $p_1$ and decreases in $p_0$ for $p_1>p_0$, it is at least $d(1-\beta\Vert\alpha)$~\cite{CoverThomas2006}. For a fixed path the chain rule gives $D(P_1\Vert P_0)=\sum_mD(p_{\theta_m}\Vert p_0)$; for a random path $P_1$ is a mixture and joint convexity gives $D(P_1\Vert P_0)\le\E\sum_mD(p_{\theta_m}\Vert p_0)$. This proves~\eqref{eq:detect_bound}. For (a), summing the hypothesis over $m\le rT$ gives $\kappa r\varepsilon^4T^5/40\,(1+O(\varepsilon^2T^2))$, and the solution $T\propto\varepsilon^{-4/5}$ has $\varepsilon T\to0$. For (b), summation by parts gives $\sum_{t_m\le T,\,m\ {\rm probe}}t_m^2=p rT^3/3+o(T^3)$. The probe records therefore give $pF_1r\varepsilon^2T^3/6\,(1+o(1)+O(\varepsilon^2T^2))$ and the in-frame records only $O(\varepsilon^4T^5)$, which vanishes at $T\propto\varepsilon^{-2/3}$. For (c), the generator $\sigma\Delta$ on the circle makes the unwrapped angle a Brownian motion with variance $2\sigma t$, and the relative entropy~\eqref{eq:kl_inframe} is periodic and bounded, so $D(p_\theta\Vert p_0)\le\kappa\theta^4/8+C\theta^6$ for all real $\theta$. With $\E\theta_t^4=12\sigma^2t^2$ and $\E\theta_t^6=120\sigma^3t^3$ the sum is $\kappa r\sigma^2T^3/2\,(1+O(\sigma T))$, and $\sigma T\to0$ at $T\propto\sigma^{-2/3}$. The families of Corollary~\ref{cor:three_records} satisfy the hypotheses of (a) and (b): in (a) with $C=0$, in (b) and (c) because the relative entropies are bounded and smooth.
\end{proof}

The bound holds for any test, whatever statistic it uses. For a known drift it is also of the right order: the log-likelihood ratio has mean equal to the accumulated relative entropy and, for small drift, variance about twice that mean, so the likelihood-ratio test succeeds at a constant multiple of the bound. In the example of Section~\ref{sec:example} it reaches equal errors of $0.05$ at about $1.3$ times the bound. Evenness therefore costs time. A frame that sees only its own channel needs a time of order $\varepsilon^{-4/5}$ to notice a drift of rate $\varepsilon$; a stream that includes probes outside the frame needs a time of order $\varepsilon^{-2/3}$, which is shorter for slow enough drift at any fixed fraction $p$ of probes.

\begin{corollary}[Detection before failure]
\label{cor:before_failure}
For Paper~II's reader with $D=2$ under spherical Brownian drift, a frozen self-model becomes worse than no memory at $t^*=\ln2/\sigma$. Let the agent compare the mean $S_T$ of the squared in-frame residuals $e_m=y_m-\rho\,v\cdot x_m$ over a window of length $T$ with the threshold that gives type~I error $\alpha$ under calibration. For every $\beta>0$ there is a constant $K$, depending on $\alpha$, $\beta$, $\rho$, $s$ and $r$, such that at $T=K\sigma^{-2/3}$ the type~II error is at most $\beta$ for all small enough $\sigma$. Together with Proposition~\ref{prop:detect}(c), the detection time is of order $\sigma^{-2/3}$, and its ratio to $t^*$ tends to zero: slow drift can be detected from in-frame records long before the frozen model fails.
\end{corollary}

\begin{proof}
Paper~II's frozen-model curve with $D=2$ gives $t^*$. Given the path, $e_m\sim\mathcal N(0,s^2+2\rho^2(1-\cos\theta_m))$ independently, and under calibration $nS_T/s^2$ is $\chi^2_n$ with $n=rT$, so the threshold is $s^2(1+z_\alpha\sqrt{2/n}\,(1+o(1)))$. Write $\theta_{uT}=\sqrt{2\sigma T}\,B_u$ with $B$ a standard Brownian motion. Then the conditional mean of $S_T-s^2$ is $2\sigma T\rho^2Z_n(1+O(\sigma T))$ with $Z_n\to Z=\int_0^1B_u^2\,du$, and its conditional fluctuation is $s^2\sqrt{2/n}\,(1+o(1))$ times a standard normal variable. At $T=K\sigma^{-2/3}$ the ratio of the shift to the null scale is $\sqrt2\,\rho^2\sqrt r\,K^{3/2}Z/s^2$, independent of $\sigma$, so the type~II error tends to $P\bigl(N+\sqrt2\rho^2\sqrt rK^{3/2}Z/s^2\le z_\alpha\bigr)$ with $N$ standard normal and independent of $Z$. Since $Z>0$ almost surely, this tends to zero as $K\to\infty$.
\end{proof}

For a general predetermined probe schedule, let $\chi_m$ indicate a probe
at $t_m=m/r$. Under the linear drift of
Proposition~\ref{prop:detect}(b), the leading term of the probe entropy upper bound is
\[
 \frac{F_1\varepsilon^2}{2}\sum_{m\le rT}\chi_m t_m^2,
\]
with remainder $O(\varepsilon^4 rT^5)$ under its global entropy bounds.
The prefix-density assumption makes this sum $p rT^3/3+o(T^3)$.
If only the horizon count $\sum\chi_m=p rT+o(T)$ is specified,
placing all probes first gives $p^3rT^3/3$, while placing them last gives
$[1-(1-p)^3]rT^3/3$, to leading order. These are the sharp extremes
because the weights $t_m^2$ increase. A probe fraction alone therefore
does not determine the detection constant. The placement is a genuine
use of the access budget, even though all three schedules have the
same $\varepsilon^{-2/3}$ scaling of the necessary detection time at fixed $p>0$.

Detection is not correction. An even record can signal that the frame has moved but not in which direction (Proposition~\ref{prop:first_order_nonidentifiability}). One off-frame read or one reference read supplies the sign and a first-order signal at the same time. This is the division of labour in Paper~II's calibratable frame and in Section~\ref{sec:loop_world}: most reads serve the task through the frame, and a small budget of reads outside it keeps the frame correct.

% ============================================================
\section{Estimating the mismatch with a prior self-model}
\label{sec:SRCR}

For an agent that retains a model of its instrument, earlier calibration records supply a prior $\pi$ on the drift parameter $\delta\vartheta$. Its Fisher information measures how firmly that estimate is held. Paper~II's foreground uses such a retained model; Paper~I's implicit protocol distribution alone does not require one.

\begin{definition}[Rigidity of a self-model]
\label{def:ego_rigidity}
The rigidity of a prior self-model $\pi$ is its Fisher information
\begin{equation}
\label{eq:ego_rigidity}
\mathcal I_{\mathrm{self}}:=\int\bigl(\partial_{\delta\vartheta}\log\pi(\delta\vartheta)\bigr)^2\,\pi(\delta\vartheta)\,d(\delta\vartheta).
\end{equation}
\end{definition}

\begin{theorem}[Van Trees bound with a prior self-model]
\label{thm:SRCR}
Let $\delta\vartheta$ lie in a compact interval $[a,b]$, let the prior $\pi$ be $C^1$ with finite rigidity and vanish at $a$ and $b$, and let the record consist of $n_{\mathrm{eff}}$ conditionally independent samples (or a joint experiment with that exact information), each with Fisher information $\mathcal I_1(\delta\vartheta)$ under a regular likelihood. Then every estimator $\delta\hat\vartheta$ based on the record satisfies
\begin{equation}
\label{eq:SRCR}
\E\bigl[(\delta\hat\vartheta-\delta\vartheta)^2\bigr]\ \ge\ \frac1{n_{\mathrm{eff}}\,\E_\pi[\mathcal I_1(\delta\vartheta)]+\mathcal I_{\mathrm{self}}}.
\end{equation}
\end{theorem}

\begin{proof}
The van Trees inequality~\cite{vanTrees1968} bounds the Bayesian mean squared error of any estimator by the inverse of the prior-averaged Fisher information of the record plus the Fisher information of the prior; the boundary conditions remove the boundary term of its integration by parts. Conditional independence makes the record's information $n_{\mathrm{eff}}\mathcal I_1$.
\end{proof}

The prior helps when it is well placed and hinders when it is not. In the conjugate Gaussian location model, with constant per-sample precision $\bar{\mathcal I}_1$, the posterior mean is the precision-weighted average
\begin{equation}
\label{eq:optimal_estimator}
\hat\vartheta=\frac{n_{\mathrm{eff}}\bar{\mathcal I}_1\,\hat\vartheta_{\mathrm{MLE}}+\mathcal I_{\mathrm{self}}\,\vartheta_0}{n_{\mathrm{eff}}\bar{\mathcal I}_1+\mathcal I_{\mathrm{self}}}
\end{equation}
of the maximum-likelihood estimate and the prior centre $\vartheta_0$. A rigid self-model centred on the wrong value pulls every estimate toward it and resists recalibration; a soft one follows the noise. Without a model of the drift one can only say that an optimal rigidity exists when the loss is lower semicontinuous on a compact set of rigidities. Under a random-walk drift the optimum is exact (Corollary~\ref{cor:rigidity}).

% ============================================================
\section{The price of keeping up}
\label{sec:price}

\subsection{The drift law}
\label{subsec:driftlaw}

Let a scalar self-parameter, for instance the gain of the agent's instrument, drift as a Gaussian random walk $\theta_k=\theta_{k-1}+\omega_k$, with independent $\omega_k\sim\mathcal N(0,\nu^2)$ per use and $\theta_0$ drawn from a prior with finite Fisher information $Q_0$ ($Q_0=1/P_0$ for a Gaussian prior of variance $P_0$). In use $k$ the agent makes $\rho_k\ge0$ reference reads, each costing $w$ and carrying Fisher information at most $J$ about $\theta_k$ under the pointwise or Bayesian conditional-information budget of Lemma~\ref{lem:stopped_score}, together with any unlabelled reads, which carry expected information at most $J_0$ per use at no extra cost. It then forms an estimate $\hat\theta_k$ from all its records so far and incurs a loss. The number $\rho_k$ may be a real reading intensity, and it may depend on the records of earlier uses. Define $\bar\rho_k=\E\rho_k$, $\psi(Q)=(Q^{-1}+\nu^2)^{-1}$ and
\begin{equation}
\label{eq:pcrb}
Q_k=\psi(Q_{k-1})+J_0+J\bar\rho_k,\qquad P_k=Q_k^{-1}\qquad(k\ge1).
\end{equation}

\begin{lemma}[Predictable design and stopped reading scores]
\label{lem:stopped_score}
Fix a use $k$ with $\theta_k$ unchanged during its reads. Let
$\mathcal H_{j-1}$ contain the preceding complete acquired history,
selected designs and private coins, including earlier reads of this use,
and set $\mathcal G_{j-1}=\mathcal H_{j-1}\vee\sigma(\theta_0,\ldots,\theta_k)$.
Assume a parameter-independent, nonanticipating design rule and regular
dominated likelihoods. The next reading score $s_j$ has conditional mean
zero given $\mathcal G_{j-1}$, and put
$I_j=\E[s_j^2\mid\mathcal G_{j-1}]$.
Either the pointwise bound $I_j\le J$, or the weaker Bayesian bound
$\E[I_j\mid\mathcal H_{j-1}]\le J$ under the actual joint prior and experiment,
suffices. For a stopping count $\rho_k$ with $\E\rho_k<\infty$ and
$\{\rho_k\ge j\}\in\mathcal H_{j-1}$, the Bayesian expected reading
information is at most $J\E\rho_k$. A separately declared free record
channel contributes its joint conditional-information budget $J_0$.
\end{lemma}
\begin{proof}
The joint likelihood factors as
$p(\theta_0)\prod_k p(\theta_k\mid\theta_{k-1})
\prod_j p(z_{kj}\mid\theta_k,d_{kj})$ times design probabilities.
At the realized record the selected designs are fixed under
differentiation, so design contributes no parameter score.
The scores are martingale differences in the enlarged filtration
$\mathcal G_j$, not necessarily in $\mathcal H_j$.
Writing $A_j=\mathbf1\{\rho_k\ge j\}$, orthogonality and the tower property give
\[
 \E S_{\rho_k\wedge m}^2
 =\sum_{j\le m}\E[A_j I_j]
 =\sum_{j\le m}\E[A_j\E(I_j\mid\mathcal H_{j-1})]
 \le J\E(\rho_k\wedge m).
\]
The same argument bounds the squared $L^2$ tail by
$J\E(\rho_k-m)_+\to0$. The transition information blocks of the posterior
Cram\'er--Rao recursion are unchanged by this factorization.
The conditional-average alternative proves a Bayesian information bound,
not a separate bound at every $\theta_k$. Parameter-dependent design
or informative status records need their own scores; an unconditional
per-read average alone does not justify stopping on the record.
\end{proof}

\begin{theorem}[Drift law]
\label{thm:driftlaw}
Let $\ell:(0,\infty)\to[0,\infty)$ be a continuous function such that the expected loss in use $k$ is at least $\ell(P_k)$ for every estimator. Then every reading schedule and every estimator satisfy
\begin{equation}
\label{eq:driftlaw}
\liminf_{K\to\infty}\frac1K\sum_{k=1}^K\bigl(w\bar\rho_k+\E[\mathrm{loss}_k]\bigr)\ \ge\ \inf_{0<P\le P_{\mathrm{free}}}h(P),\qquad
h(P)=\frac wJ\Bigl(\frac{\nu^2}{P(P+\nu^2)}-J_0\Bigr)+\ell(P),
\end{equation}
where $P_{\mathrm{free}}=1/Q_{\mathrm{free}}$ and $Q_{\mathrm{free}}$ is the positive fixed point of $Q=\psi(Q)+J_0$ ($P_{\mathrm{free}}=\infty$ if $J_0=0$). For the quadratic loss $\E(\hat\theta_k-\theta_k)^2$ times a constant $c$ the hypothesis holds with $\ell(P)=cP$. If the reads are Gaussian, the loss of the posterior-optimal decoder is
exactly $\ell(P_k)$, and arbitrary nonnegative observation precision is an
allowed control, and the initial prior is Gaussian, a constant precision attaining a minimizing $P^*$ attains
the bound. When reads must be counted in whole units, the same expression
is a lower bound from a continuous-precision relaxation; equality requires
an admissible schedule.
\end{theorem}

\begin{proof}
\emph{Information bound.} For a random walk with increment variance $\nu^2$ the posterior Cram\'er--Rao recursion of Tichavsk\'y, Muravchik and Nehorai~\cite{Tichavsky1998} has $D_{11}=-D_{12}=-D_{21}=\nu^{-2}$ and $D_{22}=\nu^{-2}+I_k$, where $I_k$ is the expected Fisher information of the records of use $k$ about $\theta_k$, and it reads $\mathcal J_k=D_{22}-D_{21}(\mathcal J_{k-1}+D_{11})^{-1}D_{12}=\psi(\mathcal J_{k-1})+I_k$, with $\E(\hat\theta_k-\theta_k)^2\ge\mathcal J_k^{-1}$. Lemma~\ref{lem:stopped_score} supplies $I_k\le J_0+J\bar\rho_k$ for predictable designs, including finite-mean within-use stopping counts. Whole-use Gaussian precision has its declared conditional information, without a discrete stopping argument. Since $\psi$ is increasing, induction gives $\mathcal J_k\le Q_k$, so the mean squared error is at least $P_k$, and the quadratic loss is at least $cP_k$.

\emph{Telescoping.} By~\eqref{eq:pcrb}, $J\bar\rho_k=Q_k-\psi(Q_{k-1})-J_0$. Summing,
\[
\sum_{k=1}^K\bigl(Q_k-\psi(Q_{k-1})\bigr)=\sum_{k=1}^{K-1}\bigl(Q_k-\psi(Q_k)\bigr)+Q_K-\psi(Q_0)\ \ge\ \sum_{k=1}^{K-1}\bigl(Q_k-\psi(Q_k)\bigr)-\psi(Q_0),
\]
and $Q-\psi(Q)=\nu^2/(P(P+\nu^2))$ for $P=1/Q$. Hence
\[
\sum_{k=1}^K\bigl(w\bar\rho_k+\E[\mathrm{loss}_k]\bigr)\ \ge\ \sum_{k=1}^{K-1}h(P_k)+\ell(P_K)-\frac wJ\bigl(\psi(Q_0)+J_0\bigr)\ \ge\ \sum_{k=1}^{K-1}h(P_k)-\frac wJ\bigl(\psi(Q_0)+J_0\bigr).
\]
Each term of the sum is paired with a single $P_k$: no convexity is needed.

\emph{Range.} Since $\bar\rho_k\ge0$, $Q_k\ge f(Q_{k-1})$ with $f(Q)=\psi(Q)+J_0$ increasing, so $Q_k\ge f^k(0)$, and $f^k(0)$ increases to $Q_{\mathrm{free}}$. For every $\eta>0$, $P_k\le(1+\eta)P_{\mathrm{free}}$ for all large $k$. Dividing by $K$, letting $K\to\infty$ and then $\eta\to0$, which is allowed because $h$ is continuous, gives~\eqref{eq:driftlaw}.

\emph{Attainment.} With a Gaussian prior and Gaussian measurements, the allowed constant precision makes $Q_k$ converge to $1/P^*$. A posterior-optimal decoder with loss exactly $\ell(P_k)$ then has average cost converging to $h(P^*)$.
\end{proof}

\begin{corollary}[Paper~I's drift law]
\label{cor:driftlaw}
Take Paper~I's quadratic loss $\tfrac12k_BTG(\hat\theta-\theta)^2$ per use and $J_0=0$, and let
\begin{equation}
\label{eq:driftlaw_explicit}
C=3\cdot4^{-2/3}\,w^{1/3}\,(k_BTG\nu)^{2/3}\,J^{-1/3}.
\end{equation}
Every schedule and estimator has long-run average reference cost plus loss at least $C-\tfrac12k_BTG\nu^2$ per use. With a Gaussian prior and reads of controllable real precision, the best stationary policy achieves at most $C$, so the least achievable average $L^*$ in that class satisfies $C-\tfrac12k_BTG\nu^2\le L^*\le C$.
The optimal reference rate is $\rho^*=\bigl(k_BTG\nu/(4w\sqrt J)\bigr)^{2/3}(1+o(1))$ as $\nu\to0$, and $C=3w\rho^*$ to leading order. The optimal estimator is a persistent memory: the steady-state Kalman filter updates its estimate with the constant weight $\alpha^*=P^*J\rho^*=\nu\sqrt{\rho^*J}\,(1+o(1))$, which is the memory of Proposition~\ref{prop:memory_drift} with that weight.
\end{corollary}

\begin{proof}
With $c=\tfrac12k_BTG$ and $Q=P+\nu^2$, the inequalities $\nu^2/Q^2\le\nu^2/(P(P+\nu^2))\le\nu^2/P^2$ give $h(P)\ge(w/J)\nu^2/Q^2+cQ-c\nu^2$ and $h(P)\le(w/J)\nu^2/P^2+cP$. Both $\min_{Q>0}[(w/J)\nu^2/Q^2+cQ]$ and $\min_{P>0}[(w/J)\nu^2/P^2+cP]$ equal $3\cdot2^{-2/3}c^{2/3}(w\nu^2/J)^{1/3}=C$. Theorem~\ref{thm:driftlaw} gives the lower bound and, with Gaussian reads, the upper bound. At the stationary point $\nu^2/(P^*(P^*+\nu^2))=J\rho^*$ and $P^*=\nu/\sqrt{\rho^*J}\,(1+o(1))$; minimizing $w\rho+c\nu/\sqrt{\rho J}$ gives $\rho^*$ and the value $3w\rho^*$. The stationary Kalman gain is the posterior variance times the precision of the use's reads, $P^*J\rho^*$.
\end{proof}

Paper~I conjectured this law with an unspecified universal constant, the rate $\rho^*$ and the value $3w\rho^*$~\cite{Liu2026TDOME_I}. Corollary~\ref{cor:driftlaw} proves it, with the constant $3\cdot4^{-2/3}\approx1.191$, for every schedule including time-varying and adaptive ones in its scalar Gaussian-transition model, for any read family whose Fisher information is bounded by $J$, and for the quadratic loss defined by Paper~I's metric $G$. For Paper~I's exact divergence loss the statement holds to leading order in the calibration error; a global version holds with $G$ replaced by any $G'>0$ such that the loss is at least $\tfrac12k_BTG'(\hat\theta-\theta)^2$ on the explored range, under a uniform quadratic separation of posteriors on that range. Injectivity alone does not guarantee a positive quadratic minorant.

\begin{corollary}[The relative-entropy floor of Consciousness~III for all policies]
\label{cor:c3}
In the Gaussian self-model of Consciousness~III~\cite{Liu2026CIII}, let the initial prior be Gaussian and the self-state $\mathcal N(\theta_t,v)$ drift with increment variance $q$, with readings $\theta_t+\mathcal N(0,r)$ costing $w$ each. Every predictable observation-precision schedule and decoder has long-run average relative-entropy loss plus reading cost at least
\[
\min_{P>0}\Bigl[\frac{wrq}{P(P+q)}+\tfrac12\ln\Bigl(1+\frac Pv\Bigr)\Bigr],
\]
For a Gaussian initial prior and controllable real observation precision at linear cost, the constant-rate Kalman schedule attains it. Integer reading counts give the corresponding relaxation bound. The stationary constant-rate optimum computed in Consciousness~III is therefore optimal among all policies.
\end{corollary}

\begin{proof}
Let $\mathcal H_k$ contain the full readings, chosen precisions, and private
randomness through use $k$. A precision chosen from $\mathcal H_{k-1}$ does
not contribute an extra parameter-dependent likelihood factor. Gaussian
conjugacy therefore gives, path by path,
\[
 P_k^{-1}=(P_{k-1}+q)^{-1}+\rho_k/r,
 \qquad \theta_k\mid\mathcal H_k\sim\mathcal N(m_k,P_k).
\]
The posterior predictive law is $\mathcal N(m_k,v+P_k)$, with conditional
optimal relative-entropy loss $\ell(P_k)=\tfrac12\ln(1+P_k/v)$.
Put $Q_k=P_k^{-1}$ and $\psi(Q)=(Q^{-1}+q)^{-1}$.
For every realized observation schedule,
\[
 \sum_{k=1}^K[w\rho_k+\ell(P_k)]
 =\sum_{k=1}^{K-1}
 \left[\frac{wrq}{P_k(P_k+q)}+\ell(P_k)\right]
 +\ell(P_K)+wr[Q_K-\psi(Q_0)].
\]
The terminal terms are at least $-wr\psi(Q_0)$.
Take expectations and divide by $K$. A constant allowed precision at the
minimizer attains the resulting bound.
\end{proof}

\subsection{What the floor means}
\label{subsec:floor_meaning}

\begin{remark}[Local thermal comparison]
\label{cor:thermo}
At a fixed interior gain $g$, the Bernoulli reference information and Paper~I's local loss metric are $J=G=1/(1-g^2)$. Substitution into the exact Gaussian comparison yields
\begin{equation}\label{eq:thermo_floor}
3\cdot4^{-2/3}(\ln2)^{1/3}\frac{k_BT\nu^{2/3}}{(1-g^2)^{1/3}}
-\frac{k_BT\nu^2}{2(1-g^2)}.
\end{equation}
This expression gives the local quadratic comparison at thermally prepared reference cost $k_BT\ln2$. A Gaussian random walk leaves the physical gain interval almost surely, so this substitution alone does not prove an infinite-time lower bound for the Bernoulli instrument. A bounded-drift theorem or a uniform finite-window limit is needed for that conclusion. The Gaussian theorem itself retains its all-policy strength.
An alternative coordinate is $p=\sin^2\phi$, whose Bernoulli Fisher
information is $4$ where $0<p<1$. This alone cannot supply a global
quadratic KL minorant on the circle: $\phi$ and $\pi-\phi$ give identical
laws while their circle distance can be nonzero. A bounded-instrument
extension must use an identifiable quotient or chart, establish its
loss separation and treat its actual transition law; the Gaussian
information recursion is not coordinate-invariant under that change.
\end{remark}

Paper~I's amortization rule says that the overhead of a static calibration falls as $U^{-1/2}$ with the number $U$ of uses, and so can be made as small as desired. The exact Gaussian comparison removes that freedom: with no free information its price per use has order $k_BT\nu^{2/3}$. For $g=0.55$ the local Gaussian comparison value is $0.012\,k_BT$ per use at $\nu=10^{-3}$, where the calibration error at the optimum is about $0.1$ in $g$ and the leading-order statement is accurate to a few per cent; the quadratic-loss value at $\nu=10^{-2}$ is $0.055\,k_BT$, but there the error is about $0.23$ and the divergence loss is about $7\%$ below its quadratic approximation (Section~\ref{sec:example}). A record read through the same channel is worth $0.160\,k_BT$ per use in Paper~I's example. For a bounded physical gain, a uniform information bound and a quadratic loss minorant must be paired with the actual transition law; clipping the Gaussian walk changes its information recursion. With the quadratic loss the two would be equal near $\nu=0.05$, where calibration errors are no longer small, so that crossing is only indicative. A physical persistence calculation must combine the actual record reset and the cost of its admissible calibration policy; Section~\ref{subsec:J_viability} supplies such a bounded example.

\begin{corollary}[Optimal rigidity under drift]
\label{cor:rigidity}
In the optimal stationary policy of Corollary~\ref{cor:driftlaw}, the information that the self-model carries into each use from earlier uses is $\mathcal I_{\mathrm{self}}^*=1/(P^*+\nu^2)$, the information gathered in the use is $J\rho^*$, and their ratio is $(1-\alpha^*)/\alpha^*=(\nu\sqrt{\rho^*J})^{-1}(1+o(1))$.
\end{corollary}

\begin{proof}
At stationarity $Q^*=\psi(Q^*)+J\rho^*$ with $\psi(Q^*)=1/(P^*+\nu^2)$, and the Kalman weight is $\alpha^*=J\rho^*/Q^*$.
\end{proof}

A more rigid self-model would ignore real drift; a softer one would spend references on noise. The trade-off that Section~\ref{sec:SRCR} could only assert exists has, under drift, an exact optimum set by the drift speed and the price of a reference.

\begin{remark}[What the drift law does not cover]
\label{rem:driftlaw_scope}
(i) \emph{Free information.} If unlabelled reads inform the parameter ($J_0>0$), the floor is lower, and when $\nu/\sqrt{J_0}$ is small enough the agent can track without references at an error of about $\nu/\sqrt{J_0}$; the $\nu^{2/3}$ law then no longer describes the optimum. In Paper~I's setting with prior $1/2$, unlabelled reads inform the bias and not the gain. In a world that the agent changes, an unlabelled read has mean $gx+b$ with the world state $x$ unknown, so world reads alone do not separate the instrument from the world, which is Paper~II's blind spot in another form. Section~\ref{sec:loop_world} uses references for this reason; Section~\ref{sec:autonomous} shows that a closed loop supplies free information about some invariant self-parameters in its stationary nondegenerate model, while the symmetry obstructs absolute sign identification.
(ii) \emph{Linear drift.} Under the linear drift of Proposition~\ref{prop:memory_drift}, a memory with a constant weight has mean squared error $\nu^2(1-\alpha)^2/\alpha^2+\alpha\sigma^2/(2-\alpha)$; optimizing the weight and the number of reads gives a cost proportional to $\nu^{2/5}$, worse than $\nu^{2/3}$. A memory that also estimates the drift rate drives the error to zero under exactly linear drift, so linear drift has no positive floor.
(iii) \emph{Vector parameters.} For a self-model with several drifting parameters the same telescoping, applied to traces, gives a lower bound by relaxation; whether it is attained when one reference informs several parameters at once is open.
\end{remark}

\begin{conjecture}[Smoother drift]
\label{conj:smooth}
Let the self-parameter be an $m$-fold integrated random walk, so that its $m$-th differences are independent $\mathcal N(0,\nu^2)$ variables, and let the other inputs be as in Theorem~\ref{thm:driftlaw} with quadratic loss. Then every policy pays at least $c_m\,w^{(2m-1)/(4m-1)}\,\nu^{2/(4m-1)}$ per use, where $c_m>0$ depends only on $m$, $J$ and the loss constant, and the stationary Kalman filter with a constant reference rate attains this order.
\end{conjecture}

The stationary part is a steady-state computation: for an $m$-fold integrator observed with noise variance $R$ the stationary position error scales as $\nu^{1/m}R^{(2m-1)/(2m)}$, and optimizing $R=1/(\rho J)$ gives the exponent $2/(4m-1)$, which is $2/3$, $2/7$ and $2/11$ for $m=1,2,3$; a numerical solution of the Riccati equation reproduces these slopes. The open part is the lower bound for time-varying and adaptive schedules. The telescoping of Theorem~\ref{thm:driftlaw} does not extend directly, because the information recursion is matrix-valued and the reads inform only the first component. The conjecture would fail if a schedule that alternates bursts of reads with silence could exploit the smoothness to beat the stationary exponent.

% ============================================================
\section{The calibration loop}
\label{sec:loop}

\subsection{The natural-gradient update}
\label{subsec:update_law}

A loop that keeps the self-model current can move $\vartheta$ along the steepest descent of a frame loss in the geometry of its own records,
\begin{equation}
\label{eq:update_law}
\dot\vartheta=-\eta\,g(\vartheta)^{-1}\nabla_\vartheta L(\vartheta,t),
\end{equation}
where $\eta>0$ is the adaptation rate, $g$ is the Fisher metric of the records the loop uses, and $L(\vartheta,t)$ is the risk of acting at time $t$ with self-model $\vartheta$. For Paper~II's rank-one frame $v$, $L(v,t)=s^2+\rho^2\bigl(1-(v\cdot u_t)^2\bigr)$, which is smallest for the frame aligned with $u_t$. The update is invariant under reparametrization of $\Theta$; this matters because no coordinate system on a space of frames or instrument models is preferred. It is a definite algorithm. An implementation must supply the likelihood, evaluate or approximate its score and realize the update law; accumulating an unlabelled second moment of residuals is not in general a Fisher computation.

The metric must be nondegenerate where the loop operates, and records read through the frame fail exactly there.

\begin{proposition}[The in-frame metric degenerates at calibration]
\label{prop:degenerate}
For Paper~II's two-channel reader with the frame at angle $\theta$ to the coupling, the metric of in-frame records is $F(\theta)$ of~\eqref{eq:fisher_even_II}, and the natural gradient of $L=s^2+\rho^2\sin^2\theta$ in this metric is
\[
\frac{\partial_\theta L}{F(\theta)}=\frac{\rho^2\sin2\theta}{F(\theta)}=\frac{2\rho^2}{\kappa\,\theta}\bigl(1+O(\theta^2)\bigr),
\]
which diverges at $\theta=0$. If a fraction $p\in(0,1]$ of the records are off-frame probes, the metric $(1-p)F(\theta)+pF_\perp(\theta)$ is at least $p\rho^2\cos^2\theta/(\rho^2+s^2)$, and for $|\theta|<\pi/2$ its absolute natural gradient is bounded by $2(\rho^2+s^2)|\tan\theta|/p$.
\end{proposition}

\begin{proof}
The first statement follows from $F(\theta)=\kappa\theta^2+O(\theta^4)$ (Corollary~\ref{cor:three_records}). In~\eqref{eq:fisher_probe} the factor $\rho^2\sin^2\theta+\rho^2+s^2$ is at least the denominator's base $\rho^2+s^2-\rho^2\sin^2\theta$, so $F_\perp(\theta)\ge\rho^2\cos^2\theta/(\rho^2+s^2-\rho^2\sin^2\theta)\ge\rho^2\cos^2\theta/(\rho^2+s^2)$; divide $\rho^2\sin2\theta$ by $p$ times this.
\end{proof}

A loop that uses only in-frame records therefore takes ever larger natural-gradient steps as it approaches calibration, and by Proposition~\ref{prop:first_order_nonidentifiability} its records do not even fix the sign of the step. This is the geometric form of Paper~II's blind spot, and the remedy is the same: a fixed fraction of reads outside the frame makes the metric positive at calibration, with value $pF_\perp(0)=p\rho^2/(\rho^2+s^2)$.

The local strong-convexity tracking theorem and proof are retained in
Appendix~\ref{subsec:lyapunov}. Its assumptions require a nondegenerate
metric and a smooth geodesic neighborhood; they do not repair the
in-frame degeneracy. Persistent excitation is a separate input.

\begin{definition}[Persistent excitation]
\label{def:PE_constant}
The records are persistently exciting with constant $\alpha_{\mathrm{PE}}>0$ if for some window $T$
\begin{equation}
\label{eq:PE_constant}
\frac1T\int_t^{t+T}g(\vartheta(s))\,ds\ \succeq\ \alpha_{\mathrm{PE}}\,I\qquad\text{for all }t.
\end{equation}
\end{definition}

Pointwise regularity and a positive entropy rate of the environment do not by themselves give a uniform lower bound of this kind, and no equality between $\alpha_{\mathrm{PE}}$ and $m_L$ is assumed. Proposition~\ref{prop:degenerate} is a concrete failure of persistent excitation: in-frame records near calibration have $\alpha_{\mathrm{PE}}\to0$, and probes restore it.

\subsection{The four links of the loop}
\label{subsec:four_part}

The series uses the following model architecture. It is a template, not a necessity theorem.

\begin{definition}[Four links]
\label{def:four_part_template}
\begin{enumerate}
\item[\textup{(I)}] \textbf{Records.} A world leaves readable, stable records, which an agent can store and use (HAFF; Paper~I gives their value).
\item[\textup{(II)}] \textbf{Reading and control interface.} Instruments read the world and actuators change it, each with a cost and a legal form (Q-RAIF; Foundation~III's legal feedback channels~\cite{Liu2026FIII}). The programme's running example of an internal algebra, a Clifford representation, can serve as the channel space; nothing here depends on it.
\item[\textup{(III)}] \textbf{Self-frame.} A retained model of how the task sits in the agent's own channels, formed from records and selecting later reads (Paper~II).
\item[\textup{(IV)}] \textbf{Calibration.} References, memory and reliability weighting that keep (III) correct while the instrument drifts, and that decide how much authority over the world each reading carries (this paper).
\end{enumerate}
\end{definition}

Each link has an exact result in a stated model, listed in Section~\ref{sec:discussion}. Section~\ref{sec:loop_world} joins (II)--(IV) in one finite model in which the agent's actions change what it reads next. Section~\ref{subsec:body_access} makes the body's working state a condition of access and identifies a task domain where body history is necessary for a higher harvest. The template is not shown to be necessary for persistence: a system that can read everything it needs, or whose coupling never changes, needs neither (III) nor (IV).

% ============================================================
A separate-register implementation can add reset, reference preparation
and slow-driving actuation costs without double counting.
Appendix~\ref{sec:cost} states that conditional bound and its full proof.
These implementation charges are not a universal persistence budget.

\section{A joint loop and its reference standard}
\label{sec:autonomous}

The memory--self--introspection route can be joined around a body that the observer reads and changes. Records consume memory purity; a model of the instrument selects feedback and harvest levels; calibration revises that model. We first give the ledger and its identification boundary, then a finite-block operating regime and a continuing joint filter that uses its body history to revise both its instrument model and its reference decisions. A further construction lets the body determine whether those readings and updates can occur; latched recovery, when physically available, preserves later access. A labelled operation supplies orientation relative to its dispatch map, the calibration counterpart of HAFF Paper~G's final cut~\cite{Liu2026HAFF_G}. Record purity, field preparation and timing are supplied resources in these constructions.

\subsection{The model}
\label{subsec:J_model}

\begin{definition}[Joint model]
\label{def:joint}
One cycle is one use; entropies are in nats and energies in units of $k_BT$.
\begin{enumerate}
\item[(i)] \emph{World.} A persistent \emph{body} bit $B_t=\pm1$, flipped by a bath with probability $d$ per cycle ($c=1-2d$), and a \emph{food} pointer $X_t=\pm1$, fresh each cycle with $P(X_t=+1)=\tfrac12$. Both have degenerate Hamiltonians and are diagonal in the pointer basis.
\item[(ii)] \emph{Records.} Each cycle fresh pure memories collide with each carrier through pointer-controlled couplings~\cite{Liu2026HAFF_E,Liu2026HAFF_F}: the food leaves two disjoint fragments $F$ and $G$ of overlap $r$, the body a fragment $F_B$ and a copy $G_B$. Record formation costs a diagonal carrier nothing and consumes memory purity (Paper~I, ledger of the collision model).
\item[(iii)] \emph{Instrument.} Reading a fragment through the setting that matches its record gives $P(m\mid x)=(1+\sigma g\,mx)/2$ with $0\le g\le\sqrt{1-r^2}$, the Helstrom gain. The \emph{magnitude} $g$ follows a specified bounded transition law; the unbounded Gaussian model of Section~\ref{sec:price} is an exact comparison and requires a separate bridge to this physical instrument; the \emph{sign} $\sigma=\pm1$ may invert with probability $\eta$ per use; the \emph{setting} is one of $D$, and a read with the wrong setting returns a fair coin. The same instrument reads food and body. (This is the gain--bias channel of Section~\ref{subsec:reference_fisher} with $b=0$, and the setting is a discrete restriction of Paper~II's direction family.)
For uniform settings, the $D$ choices are coordinate axes
$\{e_1,\ldots,e_D\}\subset S^{D-1}$; exact identification removes
$\ln D$ nats of setting uncertainty. This is not Paper~II's continuous spherical alignment
price, and the present model supplies these discrete settings rather
than deriving their formation.
\item[(iv)] \emph{Actions.} The flip of Q-RAIF Paper~B, $U=e^{-i\lambda XQ}$, which with drive occupation $N$ acts as $s\,X\rho X+(1-s)\rho$, $s=1-\epsilon_N$~\cite{Liu2026QRAIF_B}, and costs $\theta_\omega(N+a\lambda^2)$ per opportunity, $\theta_\omega=\hbar\bar\omega/k_BT$, $a=1$ if switched on. On the food: flip if $M=-1$, then harvest with fixed levels $E_\pm=-\ln\hat q(\pm)$ set from the agent's estimate, by the level protocol of Paper~I. On the body: with a gate probability $p_{\mathrm g}$ known to the agent, flip if the reading is $-1$.
\item[(v)] \emph{Memory and calibration.} The record register $M$ is reset against the copy $G$, as in the cycle bound of Paper~I, or against blank qubits. References are fragments of pointers prepared in a known state and read with the same instrument. A \emph{self-label} is the recorded quench work of the harvest, which equals the energy of the level that the flipped food occupies and so records the food value relative to the agent's own level assignment.
\end{enumerate}
\end{definition}

With an ideal actuator, no drift, the matched setting and no body, the food cycle is Paper~I's read--harvest--reset cycle; the body alone with a known instrument is the persistence task of Q-RAIF Paper~B; the instrument alone is the drifting channel of Section~\ref{sec:price}. After the flip the food has $P(X'=+1)=q=(1+s\sigma g)/2$, and the harvest returns $\ln2-H_2(q)-d(q\Vert\hat q)$, with $H_2(x)=-x\ln x-(1-x)\ln(1-x)$, by Paper~I's lemma on extraction from a qubit, so $N\to\infty$ recovers the value $I(X{:}M)$ of Paper~I.

\subsection{The ledger}
\label{subsec:J_ledger}

Every read leaves an outcome that the agent must eventually erase.

\begin{proposition}[Reading tax]
\label{prop:tax}
Let a classical register $R$ with finite alphabet be updated once per cycle by a deterministic map $(R_t,Y_t)\mapsto R_{t+1}$, after which the outcome register is returned to a standard state, all registers having degenerate Hamiltonians and the update coupling them to a thermal bath with couplings that vanish at the endpoints and no other system correlated with $(R_t,Y_t)$. Then $\beta\sum_{t\le K}W_t\ge\sum_{t\le K}H(Y_t\mid R_t)-\ln|R|$. Under these isolation assumptions, every reference read costs at least $\E\,H_2((1+g)/2)$ per read in the long run, whatever its preparation costs.
\end{proposition}

\begin{proof}
As in the cycle bound of Paper~I, unitarity and $D(\rho_B'\Vert\tau_B)\ge0$ give $\beta W_t\ge H(R_t,Y_t)-H(R_{t+1})=H(Y_t\mid R_t)+H(R_t)-H(R_{t+1})$; the last two terms telescope to at least $-\ln|R|$. Condition also on the sign, known input and register history. Fresh sensing noise gives $H(Y_t\mid R_t)\ge\E H_2((1+g_t)/2)$. A retained correlated outcome copy or newly consumed pure workspace must instead enter the reset ledger.
\end{proof}

The tax distinguishes a free reference input from a free cyclic processor. Under Proposition~\ref{prop:tax}'s isolation assumptions, an outcome stream still has a reset cost even when reference preparation is free. In a valid scalar Gaussian comparison this can supply the positive reading price $w$ and preserve the $\nu^{2/3}$ order. For the local Bernoulli comparison, adding preparation and separate outcome erasure changes $w$ from $\ln2$ to $\ln2+H_2((1+g)/2)$; the leading comparison rises by a factor $1.21$ at $g=0.55$. The bounded physical transition and joint likelihood must still be established.

\begin{proposition}[Purity available to a reset]
\label{prop:purity_copy}
Suppose the reset may couple $M$, a copy $G$ whose marginal is retained,
and $n_B$ initially pure, independent blank qubits to a thermal bath.
All registers have degenerate endpoint Hamiltonians. Then
\[
 \beta W_{\rm reset}\ge S(M\mid E_G)-n_B\ln2.
\]
Consequently, for a uniform food pointer,
\[
 \beta\,\mathrm{net}
 \le \chi(M{:}E_G)-H(M\mid X)+\Pi_B
 \le \Pi_G+\Pi_B-H(M\mid X),\qquad \Pi_B=n_B\ln2,
\]
provided the record model gives $\chi(X{:}E_G)\le\Pi_G$ and
$M-X-E_G$ is the stated conditional product model.
If the allowed reset is additionally required to consume nonnegative work,
the first upper bound can also be capped by $I(X{:}M)$.
\end{proposition}
\begin{proof}
The entropy balance for reset gives the first inequality. Subtract it from
the harvest bound $\beta W_{\rm harvest}\le I(X{:}M)$, and use
$S(M\mid E_G)=H(M)-\chi(M{:}E_G)$ and
$I(X{:}M)-H(M)=-H(M\mid X)$.
Data processing gives $\chi(M{:}E_G)\le\chi(X{:}E_G)\le\Pi_G$.
Pure blanks can themselves fund work during reset; only the added
nonnegative-reset-work condition permits the cap by the harvest alone.
\end{proof}

In the uniform-pointer task, the second copy and pure blanks are resources for reset. Preparing all consumed purity from the agent's own thermal resources removes that contribution to profit. The five-memory copy of Paper~I supplies $\chi(M{:}E_G)=0.5231$ from $5\ln2=3.466$ nats of initial purity, while a blank qubit supplies $\ln2$ of entropy capacity. The copy also serves as a readable record.

\begin{theorem}[Ledger of the joint model]
\label{thm:ledger}
Use the food feedback of Definition~\ref{def:joint}, reset against the declared copies with no additional pure blanks, and account separately for every remaining processor reset and reference preparation. Let $\mathcal C_t$ denote the classical information retained and available to a reset, whose marginal the reset preserves. Condition each register's state on that information. With separate baths and no shared erasure accounting, the long-run average net free energy per cycle satisfies
\begin{equation}
\label{eq:ledger}
\beta\,\overline{\mathrm{net}}\le\E\bigl[\ln2-H_2(q_t)\bigr]-\E\bigl[S(M_t\mid E_G,\mathcal C_t)\bigr]-\E\bigl[\ell_t+w\rho_t\bigr]-\bar\tau-\theta_\omega\E\bigl[N_t+a_t\lambda_t^2\bigr]-U_{\mathrm{body}},
\end{equation}
where $\ell_t=d(q_t\Vert\hat q_t)$ is the harvest mismatch, $w\ge\ln2$ the reference preparation cost, $\bar\tau$ the cost of the remaining outcome and model-register resets, and $U_{\mathrm{body}}\ge\E[S(Y_t\mid E_{G_B},\mathcal C_t)]+\theta_\omega\E[N_B+a\lambda^2]$ the body upkeep. Without retained side information the conditional entropy reduces to $H(M_t)-\chi(M_t{:}E_G)$. Proposition~\ref{prop:tax} applies to an isolated processor among the remaining resets. To apply Theorem~\ref{thm:driftlaw}, the complete likelihood, transition law and loss must meet its hypotheses. For constant known $s$ and fixed sign, the local loop-gain coordinate $\vartheta=s\sigma g$ has reference information $1/(s^2(1-g^2))$; a marginal self-label has information $1/(1-\vartheta^2)$. Retaining the food reading changes the label likelihood and requires its full conditional information.
\end{theorem}

\begin{proof}
Apply Paper~I's extraction lemma and conditional reset bound to each state conditioned on $\mathcal C_t$, then average. Book the other processor resets, reference preparation and Q-RAIF~B's drive and switching work once each. The displayed Bernoulli informations follow by differentiating their likelihoods. Local quadratic harvest loss alone does not supply the global loss premise of Theorem~\ref{thm:driftlaw}. Under the separate-register assumptions the actual costs add to the displayed bound.
\end{proof}

Beyond leading order the calibration loss has a closed form. An agent with a belief of variance $P$ about its gain sets its levels to $\E q$, and its expected harvest is $\ln2-H_2(\E[q\mid\text{records}])$: the agent earns what its belief allows, and the loss relative to a known gain is the Jensen gap $H_2(\E q)-\E H_2(q)$, which tends to $Ps^2/(2(1-s^2g^2))$ as $P\to0$ and saturates for broad beliefs. The harvest is increasing in $s$, for a fixed gain belief. Its Jensen-gap mismatch penalty can also increase with $s$, so total harvest and mismatch penalty have different monotonicities.

\subsection{What acting can reveal}
\label{subsec:J_gauge}

Let $Z_t$ be the hidden world state, $\theta$ the hidden parameters (instrument, actuators, world dynamics, record formation, initial law $\mu_\theta$), and let the agent's record be $(Y_1,A_1,\dots,Y_t,A_t)$, with readings $Y_t\sim K^{\mathrm s}_\theta(\cdot\mid Z_t)$, actions $A_t$ chosen by an arbitrary adaptive policy with private randomness, and transitions $Z_{t+1}\sim K^{\mathrm w}_\theta(\cdot\mid Z_t,A_t)$.

\begin{definition}[Loop symmetry]
\label{def:loopsym}
A group $\mathcal G$ acting on world states and parameters is a symmetry of the loop if $K^{\mathrm s}_{g\theta}(y\mid gz)=K^{\mathrm s}_\theta(y\mid z)$, $K^{\mathrm w}_{g\theta}(g\,\cdot\mid gz,a)=K^{\mathrm w}_\theta(\cdot\mid z,a)$ and $\mu_{g\theta}=g_*\mu_\theta$ for all $g\in\mathcal G$.
\end{definition}

\begin{theorem}[Acting cannot see the symmetry of its loop]
\label{thm:gauge}
If $\mathcal G$ is a loop symmetry, then for every policy the law of the agent's record is the same under $g\theta$ as under $\theta$, and $\E_{g\theta}L(Z_t)=\E_\theta L(gZ_t)$ for every loss $L$. A parameter function that $\mathcal G$ moves has no consistent estimator from the agent's records, however it acts. If $\gamma\in\mathcal G$ exchanges a binary loss, $L(\gamma z)=1-L(z)$, then $\E_\theta L(Z_t)+\E_{\gamma\theta}L(Z_t)=1$ for every policy and every $t$.
\end{theorem}

\begin{proof}
Write the joint density of hidden variables and records as $\mu_\theta(z_0)\prod_tK^{\mathrm s}_\theta(y_t\mid z_t)\,\pi_t(a_t\mid\cdot)\,K^{\mathrm w}_\theta(z_{t+1}\mid z_t,a_t)$ and substitute $z\mapsto gz$ in every hidden variable. The policy factors do not involve $z$, and the others are invariant by Definition~\ref{def:loopsym}. Integrating out the hidden variables gives equal record laws, and the same substitution gives the loss identity.
\end{proof}

In the joint model, $\gamma:(B,X,\sigma)\mapsto(-B,-X,-\sigma)$ with everything else fixed is a loop symmetry for any agent that only reads and flips: the instrument depends on $\sigma x$, the flip commutes with $x\mapsto-x$, the bath and the collisions are symmetric, and the gate coin is independent. The body loss $\mathbf 1\{B=-1\}$ is exchanged. Every such agent therefore satisfies $P_\theta(B_t=-1)+P_{\gamma\theta}(B_t=-1)=1$: if it holds its body at loss $L$ it holds it at $1-L$ in the world related to it by $\gamma$, with identical records. Exact enumeration over twenty random adaptive policies of five steps gives record laws that agree to $1.2\times10^{-16}$ in total variation; one reference read per step separates them by at least $0.30$. This is the observation of the Imprint Chain that its calibration statistic depends only on $\nu^2$~\cite{Liu2026Chain}, extended to every policy and to the task loss. For a linear loop, $x_{t+1}=ax_t+\kappa u_t+f+w_t$, $y_t=gx_t+b+e_t$, the affine maps $x\mapsto\lambda x+\mu$ with $(g,b,\kappa,f)\mapsto(g/\lambda,\,b-g\mu/\lambda,\,\lambda\kappa,\,\lambda f+\mu(1-a))$ form a loop symmetry when the initial and process-noise laws are transformed as well; acting can identify the loop gain $g\kappa$ and never $g$ or $\kappa$ separately (Kalman likelihoods under random affine maps and dithered adaptive policies agree to $6\times10^{-14}$).

\begin{proposition}[What the closed loop identifies]
\label{prop:loopmoments}
In the stationary body loop with constant gain and sign, $0<g\le1$, $d\in(0,1)\setminus\{1/2\}$, and the matched setting, gate coin $A_t$ of known probability $p_{\mathrm g}\in(0,1)$ and flip success $s$,
\begin{gather*}
\E[Y_{t+1}Y_t\mid A_t=0]=cg^2,\qquad
\E[Y_{t+2}Y_t\mid A_t=A_{t+1}=0]=c^2g^2,\\
\E[Y_{t+1}Y_t\mid A_t=1]=cg^2(1-s)+cs\,\E Y_t .
\end{gather*}
Writing the three moments as $K_0,K_2,K_1$, and $\mu=\E Y_t$, one has
\[
c=K_2/K_0,\qquad g^2=K_0^2/K_2,\qquad s=(K_1-K_0)/(c\mu-K_0).
\]
Stationarity gives $\mu=cp_{\mathrm g}sg^2/(1-c+cp_{\mathrm g}s)$, so the denominator is nonzero. Thus $c$, $g^2$ and $s$ are determined by the agent's own records, without references, and $\sigma$ is not. With a wrong setting the first and third moments vanish.
\end{proposition}

\begin{proof}
$B_{t+1}=\xi_tB_t\Phi_t$ with $\E\xi_t=c$ independent and $\Phi_t=-1$ if the flip is switched on and succeeds. For $A_t=0$, $\E[Y_{t+1}Y_t]=\sigma gc\,\E[B_tY_t]=cg^2$. For $A_t=1$, $\E[B_t\Phi_tY_t]=\sigma g+2s\,\E[B_t\mathbf 1\{Y_t=-1\}]$ and $\E[B_t\mathbf 1\{Y_t=-1\}]=(\E B_t-\sigma g)/2$; multiply by $\sigma gc$ and use $\sigma g\,\E B_t=\E Y_t$. The lag-two identity iterates the first.
\end{proof}

A frame error is thus visible in the agent's own loop as a loss of loop correlation; locating the right setting still requires trying the others, which is Paper~II's blind spot in binary form. What no record reveals is fixed by the following statement. Call an operation \emph{symmetry-breaking} if its action singles out a world state in a way that the symmetry does not respect: a reference prepared in the state $+1$, a reset to $+1$, harvest levels with the higher level on $-1$, a bath relaxing toward a known state. Each carries a label $\zeta$ saying which state it singles out, and that label belongs to the operation's dispatch map in the sense of Paper~G.

\begin{proposition}[No internal standard]
\label{prop:nostandard}
Let $\mathcal G$ act also on the labels of all symmetry-breaking operations by relabelling. If the labelled sensing, transition and initial kernels are jointly covariant, $\mathcal G$ remains a loop symmetry. No policy identifies an absolute jointly relabelled orientation. A system can form and retain relational standards.
\end{proposition}

\begin{proof}
Apply the assumed joint covariance in Theorem~\ref{thm:gauge}. Set-theoretic closure of a kernel family alone does not imply covariance.
\end{proof}

The model locates an operational reference standard in a labelled kernel: an external reference, an actuator reset or level assignment, a known environmental asymmetry, or a previously interpreted record. Jointly relabelling every such kernel preserves the record law. A system can nevertheless construct relational standards and carry them between operations. Paper~G places the final cut of internal selection at the dispatch map of the top executor~\cite{Liu2026HAFF_G}; here the corresponding calibration object is the labelled operation whose orientation the loop uses. The symmetry theorem does not derive the physical origin of that operation.

\subsection{Three self-parameters, three laws}
\label{subsec:J_laws}

The magnitude and the setting are invariant under the sign symmetry; the sign is not.

\begin{proposition}[Free loop information changes the exponent]
\label{prop:J0}
In Theorem~\ref{thm:driftlaw} let unlabelled reads carry $J_0>0$ per use at no extra cost, and let the loss be $cP$. Put $P_f=\tfrac12\bigl(-\nu^2+\sqrt{\nu^4+4\nu^2/J_0}\bigr)$. If $c\le(w/J)J_0(2P_f+\nu^2)/(P_f(P_f+\nu^2))$, which holds for $\nu\le\nu_\times=2wJ_0^{3/2}/(Jc)\,(1+o(1))$, every policy pays at least $cP_f=c\nu/\sqrt{J_0}\,(1+O(\nu\sqrt{J_0}))$ per use, and with Gaussian reads this is attained with no references.
\end{proposition}

\begin{proof}
$\phi(P)=\nu^2/(P(P+\nu^2))$ is convex and decreasing, so $h$ of Theorem~\ref{thm:driftlaw} is convex. $P_f$ solves $\phi(P)=J_0$, and $|\phi'(P_f)|=J_0(2P_f+\nu^2)/(P_f(P_f+\nu^2))$, so under the condition $h'\le0$ on $(0,P_f]$ and the infimum is $h(P_f)=cP_f$. As $\nu\to0$, $P_f=\nu/\sqrt{J_0}-\nu^2/2+O(\nu^3)$. The Kalman filter on the free reads alone has stationary variance $P_f$.
\end{proof}

The body moments establish static identifiability under Proposition~\ref{prop:loopmoments}'s nondegeneracy conditions. A joint information-rate bound for the correlated record stream is a further step: the information of one overlapping product is not that rate. Proposition~\ref{prop:J0} gives an exact Gaussian realization of the linear-in-$\nu$ price when a genuine free observation channel is supplied. For self-labels the marginal Bernoulli information is $1/(1-\vartheta^2)$; retaining the earlier food read changes the conditional likelihood and its information.

\begin{theorem}[Sign inversions]
\label{thm:inversion}
Suppose that sign inversions and body-bath flips are independent Bernoulli
events with probabilities $\eta,d\in(0,1)$ per use. The agent knows the initial
sign, has no symmetry-breaking operation, and reads only records formed in
the current cycle. Put
\[
 b_* = \max\{|1-2\eta|,|1-2d|\}.
\]
For every read-and-flip policy,
\[
 \left|P(B_t=-1)-\tfrac12\right|\le\tfrac12 b_*^t.
\]
In particular the loss converges to $1/2$. If
$\eta\le\min\{d,1-d\}$, its lower bound is
$\tfrac12[1-(1-2\eta)^t]$.
\end{theorem}
\begin{proof}
Let $a_j$ be the exclusive-or of the bath flip and the instrument inversion.
Condition on the apparent-body trajectory, the $a_j$, and all independent
reading and actuation randomness. This is more information than the agent
receives. Conditional inversion probabilities are
\[
 q_1=\frac{\eta(1-d)}{\eta(1-d)+d(1-\eta)},\qquad
 q_0=\frac{\eta d}{\eta d+(1-\eta)(1-d)}.
\]
Conditional on that information, the absolute sign bias is
$\prod_{j=1}^t|1-2q_{a_j}|$. The absolute conditional body bias has the same
value, since the apparent body is then fixed. Hence
\[
 |\E B_t|\le
 \left(\sum_{a=0}^1P(a)|1-2q_a|\right)^t
 =\bigl(|1-\eta-d|+|\eta-d|\bigr)^t=b_*^t.
\]
The $a_j$ are independent with a law unaffected by the policy. This proves
both sides of the bound. The special case follows from
$b_*=1-2\eta$ under the stated condition.
\end{proof}

The limit $\tfrac12$ is the loss of an agent that does nothing. An inversion is the image under $\gamma$ of a world kick that the loop corrected, so the records never show it. The food fares worse: after an inversion the flip pushes it the wrong way, the harvest runs at a loss, and an agent that never reads its energy store cannot tell. Reading the store is a self-label, whose standard is the agent's own level assignment. The Bayes-optimal policy, computed by backward induction over all adaptive policies for eight cycles, respects the bound in every case with $\eta\le d$ (for $g=0.9$, $d=0.1$, $s=0.95$, $\eta=0.05$: $0.307$ against $0.285$ at $t=8$). Consciousness~III shows that an observer can be confidently and undetectably wrong on a two-point fibre of its protocol~\cite{Liu2026CIII}; Theorem~\ref{thm:inversion} adds that acting does not change the protocol in the relevant sense, and that inversions make the error grow to chance. The hypothesis of fresh records matters. If a stored record of an earlier body state, made by the environment's fixed couplings, is read again later, the two reads compare $\sigma$ at the two times, so the world's stable records detect changes of the sign though not the sign. Records of Paper~F's kind~\cite{Liu2026HAFF_F} thus carry an agent's calibration across time, and the standard sits in the past.

\begin{proposition}[Reference cost under exogenous and adaptive timing]
\label{prop:sqrtlaw}
Let a reference cost $w>0$ and reveal the current sign exactly, and let body
loss cost $V>0$ per use. For $0\le r<1$, define
\[
 F_r(G)=\frac G2-\frac{1-r^G}{2(1-r)},\qquad
 C_r=\inf_{G\in\mathbb N}\frac{w+VF_r(G)}G.
\]
Reference times fixed in advance, or randomized independently of the record,
give long-run average cost at least $C_{b_*}$, where $b_*$ is defined in
Theorem~\ref{thm:inversion}. For reference times chosen from the records,
the bound $C_{r_*}$ holds with
\[
 r_* = \max_{a=0,1}|1-2q_a|.
\]
If $\eta\to0$ with $d\in(0,1)$ and $w,V$ fixed, these benchmarks satisfy
\[
 C_{b_*}=\sqrt{2wV\eta}\,(1+o(1)),\qquad
 C_{r_*}=\sqrt{2wV\eta\,
       \min\{d/(1-d),(1-d)/d\}}\,(1+o(1)).
\]
When a finite gap attains the first infimum, an ideal periodic
read-and-control policy attains it; otherwise the no-reference limit
gives its value $V/2$. It need not bound record-adaptive schedules.
\end{proposition}
\begin{proof}
With an exogenous gap, the mean loss at age $j$ after a reference is at
least $(1-b_*^j)/2$, by the proof of Theorem~\ref{thm:inversion}.
For an adaptive gap one cannot use this unconditional mean after
conditioning on the absence of another reference. Instead, the same proof
gives the conditional bound $(1-r_*^j)/2$ for every apparent-event history.
The event that a policy has not yet requested another reference is determined
by such a history and its private randomness, so this bound also holds on
each adaptive gap. Summing over complete gaps and the terminal incomplete
gap gives at least $KC_r-O(w)$ over $K$ uses. Equivalently, the piecewise
linear interpolation of $F_r$ is convex, and an exogenous mixture of gap
lengths cannot beat the best single length for the corresponding benchmark.
Writing $r=1-2a$ gives
$F_r(G)=aG(G-1)/2+O(a^2G^3)$; minimization has
$G\sim\sqrt{2w/(Va)}$ and cost $\sqrt{2wVa}(1+o(1))$.
For small $\eta$, $(1-b_*)/2=\eta$ and
$(1-r_*)/2=\eta\min\{d/(1-d),(1-d)/d\}(1+O(\eta))$.
Ideal reads reveal the apparent state, and ideal flips can hold it aligned
with the optimal sign estimate, attaining the exogenous benchmark.
\end{proof}
\begin{remark}[Record-adaptive references improve the benchmark]
Take $\eta=d=1/10$, $w=1/10$ and $V=1$. Read the apparent body exactly.
After accounting for the preceding commanded flip, request a reference
precisely when the apparent body has changed sign. Then control the body
using the last reference's sign estimate. Such events have probability
$p=2\eta(1-\eta)=9/50$. Between them the conditional inversion probability
is $q_0=1/82$. The stationary posterior sign bias is
\[
 z=p+(1-p)(1-2q_0)z=\frac9{10}.
\]
The loss is $1/20$, and references cost $9/500$, giving total cost
$17/250=0.068$. The exogenous-gap benchmark is $1/10$.
An exact finite-state computation with sensing gain and flip success both
$0.999$ gives cost $0.0687253$, so the separation persists away from the ideal
limit. The bound for adaptive timing is $C_{r_*}=0.0429971$ here.
\end{remark}



\begin{proposition}[Exact adaptive value at equal inversion and bath rates]
\label{prop:adaptive_exact}
In Proposition~\ref{prop:sqrtlaw}, take $\eta=d=1/10$, $w=1/10$ and $V=1$.
The minimum long-run average cost, with ideal apparent-body reads and ideal
flips, over all adaptive reference policies is
\[
 \rho_* = \frac{2089311}{43698622}=0.0478118280\ldots.
\]
Request a reference after an apparent innovation or after five consecutive
quiet transitions since the last reference. This policy attains the value,
which also bounds policies with less informative sensing or control.
\end{proposition}
\begin{proof}
An apparent innovation has probability $p=9/50$ and conditional sign bias
zero; a quiet transition has probability $z=41/50$ and multiplies the bias
by $u=40/41$. The oracle therefore has states $k=1,2,\ldots$ for $k$ quiet
transitions after a reference and a zero-bias state. At state $k$, continued
control has loss $l_k=(1-u^k)/2$; a reference resets the bias and costs $w$.
For threshold $K=5$, a renewal cycle has
\[
 D_K=\sum_{j=0}^{K-1}z^j,\qquad
 C_K=w+\sum_{j=1}^{K-1}z^jl_j,\qquad \rho_*=C_K/D_K.
\]
To verify global optimality, set $v_k=w$ for $k\ge5$ and set
$v_k=l_k-\rho_*+pw+zv_{k+1}$ backwards for $k=4,3,2,1$.
Exact rational arithmetic gives
\[
 -\rho_*+pw+zv_1=0,\quad v_k<w\ (1\le k<5),\quad l_5>\rho_*.
\]
Thus $v_k=\min\{w,l_k-\rho_*+pw+zv_{k+1}\}$ for every $k$;
monotonicity of $l_k$ proves the stopping inequality for all $k\ge5$.
The zero-bias state also stops, since $1/2>\rho_*$. These bounded Bellman
inequalities telescope over every adaptive, randomized or nonstationary
policy, leaving a bounded terminal term. Dividing by the horizon proves
the lower bound. Ideal readings reveal the apparent innovations, and ideal
flips attain their conditional Bayes body loss, so the threshold policy is
an admissible implementation of the oracle optimum. Less informative
records are a garbling of this oracle. A less reliable flip can be
simulated by an ideal actuator using a private coin, with its storage
and reset charged to the policy. Neither restriction improves the
oracle value.
\end{proof}

In the exogenous benchmark with $V=1$, the small-$\eta$ calculation is within $4\%$ of $\sqrt{2w\eta}$ over the tested $\eta\le10^{-4}$ range. A periodic noisy-reference policy gives excess cost within a factor $2$--$3$ of $\sqrt{2wV(1-2L_0)\eta}$, where $L_0=0.052$ is its no-inversion loss. Its fitted exponent is $0.41$--$0.43$, consistent with the growing reference-block size. Proposition~\ref{prop:adaptive_exact} settles the stated finite-rate ideal example. Proposition~\ref{prop:sqrtlaw} supplies the general adaptive lower bound; its constant need not be sharp.

\begin{table}[t]
\centering
\small
\begin{tabular}{@{}P{3.0cm}P{1.7cm}P{4.0cm}P{3.9cm}@{}}
\toprule
Self-parameter & Sign symmetry & Without any standard & With references or labels\\
\midrule
Setting (switches) & invariant & detected by the loop (Proposition~\ref{prop:loopmoments}); found by search & not needed\\[3pt]
Magnitude (random walk $\nu$) & invariant & static body moments identify gain; a valid Gaussian free channel gives price $c\nu/\sqrt{J_0}$ (Proposition~\ref{prop:J0}) & $\nu^{2/3}$ (Theorem~\ref{thm:driftlaw}); marginal self-labels avoid pointer preparation; full retained history determines their information\\[3pt]
Sign (inversion rate $\eta$) & variant & loss $\to\tfrac12$ for fresh-record read-and-flip policies (Theorem~\ref{thm:inversion}) & positive square-root lower bound for adaptive timing; exact finite-rate optimum in Proposition~\ref{prop:adaptive_exact}\\
\bottomrule
\end{tabular}
\caption{What acting buys for each self-parameter of the joint model.}
\label{tab:threelaws}
\end{table}

\subsection{Operating and continuing regimes}
A bounded two-gain block model has a supplied-resource operating margin
of $0.02396\,k_BT$ per food use with a squeezed drive, body loss
$0.08725$ and explicit purity input. Its complete ledger and comparison
are in Appendix~\ref{subsec:J_viability}.
A separate continuing eight-state filter carries instrument and body
history through every use. Under the same reference cap, adaptive
timing gives body loss $.10319$, versus $.11726$ for periodic timing.
Its transition/likelihood, two-seed check, 350-bit quantization certificate
and resource boundary are retained in Appendix~\ref{subsec:online_joint}.
The block margin cannot be imported into that continuing protocol
without its own apparatus ledger.

\section{When the body controls access}
\label{subsec:body_access}

The preceding loop receives a reading and can process it at every use, even when its body is at the loss state. A body that supports those operations changes the problem. Keep the bounded instrument, record couplings, field drive and body bath, but let $B=+1$ enable the main instrument, reference ports, model logic and food protocol. At $B=-1$ these ports are disconnected. The model register is held, not updated. A disabled reader cannot acquire a reference or ask a running estimator to repair it.

\begin{definition}[Access-gated joint loop]
\label{def:access_joint}
At each working use, read the body with likelihood $(1+y\sigma g_G)/2$, update the model, perform the food task and write a recovery command $\ell\in\{0,1\}$ to a latch. The working body is not actively flipped; the bath disables it with probability $d$. During disability the latch and model register remain unchanged. A separate standby port, if present, executes that stored command through the same field drive; it obtains no new instrument reading. The recovery probability per disabled slot is
\[
 r^{\rm rec}_\ell=d+\ell s(1-2d),\qquad 0<d<\tfrac12.
\]
The instrument continues its gain and sign transitions in all wall-clock slots. Disabled-slot fragments and enable flags leave through output streams that the reader cannot later acquire; there is no bank of delayed readings on resumption. The reader obtains no elapsed-time counter on recovery. A resumed use is another working invocation, with its prior obtained from the event kernel below.
Harvest work enters the physical ledger but is not read back as an additional self-label unless a protocol explicitly permits that read. The latch is written after all records acquired in its working invocation.
\end{definition}

An independent enable flag copies the pointer value of $B$ and controls the standby drive while the main processor is held. In the diagonal body model the repair leaves $+1$ fixed and sends $-1$ to $+1$ with probability $s$. A CPTP extension has Kraus operators
\[
 D_0=P_+,\qquad D_1=\sqrt{1-s}\,P_-,\qquad
 D_2=\sqrt{s}\,\lvert+\rangle\langle-\rvert,
 \qquad \sum_jD_j^\dagger D_j=I.
\]
It dephases the body blocks and is nonunital. The enable flag, its preparation or reset, and the body's free-energy change enter the ledger. One explicit resource model supplies a fresh pure flag per wall-clock slot and retains the used flag in a charged output stream, consuming $\ln2$ of flag purity per slot. Cyclic reuse requires a conditional reset and protection protocol. Table~\ref{tab:body_access_ledger} collects these resource declarations. With no standby, the symmetric bath still recovers the body at rate $d$.

\begin{proposition}[The kernel between working uses]
\label{prop:body_event}
Let $\mathsf K$ be the four-state gain/sign factor of $T_a$ and fix a stored command with recovery rate $r_{\rm rec}>0$. If $L$ is the number of wall-clock slots between successive working invocations, then
\begin{gather}
 P(L=1)=1-d,\qquad
 P(L=k)=d r_{\rm rec}(1-r_{\rm rec})^{k-2}\quad(k\ge2),
 \label{eq:body_gap}\\
 \E L=1+d/r_{\rm rec},\qquad
 \pi_{\rm w}=\frac{r_{\rm rec}}{r_{\rm rec}+d},
 \label{eq:working_fraction}\\
 \mathsf K_\ell=(1-d)\mathsf K+
 d r_{\rm rec}\mathsf K^2[I-(1-r_{\rm rec})\mathsf K]^{-1}.
 \label{eq:body_event_kernel}
\end{gather}
$\pi_{\rm w}$ is the long-run working fraction. A right eigenmode of $\mathsf K$ of eigenvalue $\lambda$ has event eigenvalue
\[
 f_{r_{\rm rec}}(\lambda)=(1-d)\lambda+
 \frac{d r_{\rm rec}\lambda^2}{1-(1-r_{\rm rec})\lambda}.
\]
The event kernel remains valid conditionally when the latch is chosen from the preceding working record.
\end{proposition}
\begin{proof}
The body either remains working for the next slot or enters a geometric recovery wait. This gives~\eqref{eq:body_gap} and its mean. The stationary flow balance is $\pi_{\rm w}d=(1-\pi_{\rm w})r_{\rm rec}$; renewal averaging gives the same fraction. Instrument noise is independent of the body and recovery noise conditional on the selected latch. Sum $\mathsf K^L$ with the displayed gap probabilities. Since $\|(1-r_{\rm rec})\mathsf K\|_\infty=1-r_{\rm rec}<1$, the sum is the positive resolvent in~\eqref{eq:body_event_kernel}. Applying it to an eigenvector gives $f_{r_{\rm rec}}$. No recursion is executed in disabled slots: the finite controller uses the matrix for its last stored command on the next invocation.
\end{proof}

The working body provides a relative calibration standard. With static sign and gain, its working readings are independent samples of $(1+y\sigma g)/2$; their mean identifies $\sigma g$ for $g>0$. With a fixed latch their Fisher information rate is $\pi_{\rm w}/(1-g^2)$ in the signed-gain coordinate. Under a history-dependent latch, information is counted on actual acquired uses, as in Consciousness~III~\cite{Liu2026CIII}. The transformation $(B,X,\sigma)\mapsto(-B,-X,-\sigma)$ exchanges enabled and disabled blocks and hence fails to preserve this access kernel. The orientation comes from the declared body--permission coupling. Jointly relabelling that coupling, every operation and the task preserves the absolute-orientation obstruction.

The readout must be informative: $g=0$ or wrong-setting reads give no instrument information. Continued access is a separate condition. In a hardware variant with an absorbing disabled state and conditional failure hazard at least $h>0$ on every working use, every policy has
\[
 P(N\ge n)\le(1-h)^{n-1},\qquad \E N\le h^{-1},
\]
where $N$ is its total number of working readings. Successive conditioning proves the bound, so only finitely many readings occur almost surely. If two static parameters have positive overlapping likelihoods on finite stopping records, their stopped-record laws have positive overlap and their equal-prior Bayes error is strictly positive. Waiting longer cannot identify them. The absorbing variant removes passive recovery as well as standby.

\begin{proposition}[Calibration through unobserved disability]
\label{prop:body_blackout}
Take known $0<g<1$, sign inversion probability $0<\eta<1/2$, one body reading on resumption and no additional current directional observation. Conditional on a stored recovery command, every sign estimator has average error at least
\[
 \min\{e_{\rm rec},(1-g)/2\},\qquad
 e_{\rm rec}=\frac{1-f_{r_{\rm rec}}(1-2\eta)}2>0.
\]
Increasing $r_{\rm rec}$ lowers this oracle floor and increases the working fraction.
For small $\eta$, $e_{\rm rec}=\eta/\pi_{\rm w}+O(\eta^2)$.
\end{proposition}
\begin{proof}
Grant the reader the true sign at its preceding working invocation. Across an unobserved gap $L$, its sign bias is multiplied by $(1-2\eta)^L$. Averaging gives the prior error $e_{\rm rec}$. The current body reading is a binary symmetric channel of error $e=(1-g)/2$. For prior error $q\le1/2$, its two-output Bayes error is
\[
 \min\{(1-q)(1-e),qe\}+
 \min\{(1-q)e,q(1-e)\}=\min\{q,e\}.
\]
An actual history is less informative than the oracle history. For $0<\lambda<1$, differentiation gives
\[
 \frac{\partial f_{r_{\rm rec}}(\lambda)}{\partial r_{\rm rec}}
 =\frac{d\lambda^2(1-\lambda)}{[1-(1-r_{\rm rec})\lambda]^2}>0.
\]
The working fraction is also increasing. A paid elapsed-time counter or another current reference changes the available record and therefore changes this bound.
\end{proof}

Filtering at actual working invocations uses the positive four-state
event kernels of Proposition~\ref{prop:body_event}. The certificate in
Appendix~\ref{app:filter_certificates} needs 150 ideal model bits and bounds
asymptotic total-variation error by $1.247\times10^{-7}$ for the stated
switching rates. Its precision is fixed; identifying a continuous
parameter to arbitrarily growing precision requires growing memory.


\begin{proposition}[A domain where history is necessary for the food task]
\label{prop:body_memory_value}
Fix the recovery command, a known static $0<g<1$, an unknown static sign with equal prior, one body reading per working invocation and no further directional sample before the food action. Within the joint food protocol, choose harvest levels from the body model before the fresh food reading $M$, then choose whether to flip from $M$. Without retained history, the maximal harvest is
\[
 J_1=\ln2-H_2((1+s g^2)/2).
\]
Let $a_g=(1+g)/2$, $b_g=(1-g)/2$, and retain the latest three body readings. For their number $k$ of positive results set
\[
 p_k=\tfrac12\binom3k(a_g^kb_g^{3-k}+b_g^ka_g^{3-k}),
 \quad m_k=\frac{a_g^kb_g^{3-k}-b_g^ka_g^{3-k}}
 {a_g^kb_g^{3-k}+b_g^ka_g^{3-k}}.
\]
The achievable harvest
\[
 J_3=\sum_{k=0}^3p_k[\ln2-H_2((1+s g|m_k|)/2)]
\]
is strictly larger than $J_1$. Thus a target harvest in $(J_1,J_3]$ requires history within these permissions. An already learned sign encoded in a program, latch or weight counts as retained history.
\end{proposition}
\begin{proof}
The fair food reading has no marginal sign information. A body-only record of posterior sign mean $m$ allows a food polarization of at most $s g|m|$: for arbitrary flip probabilities depending on $M$, its magnitude is $s g|m|$ times the difference of two probabilities, whose absolute value is at most one. Choosing the appropriate orientation attains that endpoint, and matched levels attain the stated harvest. Without history $m=gY$. Three-sample Bayes gives $p_k,m_k$. The function $F(m)=\ln2-H_2((1+s gm)/2)$ is strictly convex, with $F''(m)=(s g)^2/[1-(s gm)^2]>0$. Posterior means form a martingale, and two extra informative noisy samples make its conditional value nonconstant. Strict Jensen therefore gives $\E F(m_3)>\E F(m_1)$. A three-bit sliding record acquires no additional samples and remains held during disability.
\end{proof}

For $g=3/5$ and $s=49/50$, exact finite body--food--drive path enumeration gives $J_1=0.0635939$ and $J_3=0.1202672$. If levels may instead depend on the current food reading, the posterior-matched protocol family of Paper~I replaces $s g$ by $g$ in both formulas and retains the strict history gap.

Performance is averaged over wall-clock slots. The fixed-latch harvest gain is $\pi_{\rm w}(J_3-J_1)$ before the extra memory cost. The drive of Section~\ref{subsec:online_joint} changes the working fraction from $1/2$ to about $0.949083$ when standby is enabled. Body readings at known $B=+1$ contain fresh instrument noise uncorrelated with the body copy. Table~\ref{tab:body_access_ledger} therefore charges their accessible reset information separately, together with holding and processing. The finite-block operating margin uses a different body protocol.

\begin{table}[ht]
\centering\small
\begin{tabular}{@{}P{3.2cm}P{4.6cm}P{4.7cm}@{}}
\toprule
Resource & Declared supply or charge & Availability and accounting\\
\midrule
Food and body fragments & $4n\ln2$ purity nats per wall slot if continuously allocated & Acquired and reset only on working uses; disabled outputs are not banked\\[3pt]
Standby enable flag & $\ln2$ supplied purity nats per wall slot in the fresh-flag implementation & Used outputs leave the reader; cyclic reuse needs a reset/protection ledger\\[3pt]
Two drive packets & $2\theta_\omega N$ per wall slot & Unused prepared packets remain charged\\[3pt]
Drive switching & $\theta_\omega\lambda^2[\ell(1-\pi_{\rm w})+\pi_{\rm w}/2]$ & Recovery on disabled slots; food on working slots\\[3pt]
Body outcomes and model & Conditional erasure, update and reliable-hold costs & No processing during disability; the body copy does not erase fresh sensing noise\\
\bottomrule
\end{tabular}
\caption{Resource declarations for access-gated comparisons. Purity is an input, not a proved total energy cost. Preparation from thermal resources and any remaining apparatus work must also be charged.}
\label{tab:body_access_ledger}
\end{table}

For a fully known instrument and fixed active pipeline, let $v$ denote net value per working use after its active costs, and put $\kappa=\theta_\omega\lambda^2$. If all base supplies, including fresh flags, are allocated equally in wall time, the two fixed latches have residual values
\[
 U_\ell=\pi_{{\rm w},\ell}v-(1-\pi_{{\rm w},\ell})\kappa\ell-C_{\rm base},
 \qquad U_1-U_0=\frac{(r_1^{\rm rec}-d)v-2d\kappa}{2(r_1^{\rm rec}+d)}.
\]
This identifies the priced gain from making future calls possible. Flag costs that depend on the latch must be added, not hidden in a common base. For an uncertain instrument, the future task value and the evidence lost while disabled must also enter the control problem. The known-parameter price threshold motivates a stronger history-dependent hypothesis.

\begin{proposition}[Equal sign confidence, opposite paid repairs]
\label{prop:repair_witness}
There is a static, clocked version of Definition~\ref{def:access_joint} in
which two legally acquired histories at the same working slot have equal
sign confidence and opposite optimal repair decisions. Take $d=1/20$,
$s=49/50$, $H=12$, and a uniform prior on the four static states
$(\sigma,g)\in\{\pm1\}\times\{3/20,g_+\}$. At slot $t=10$, compare the
history with positive readings at slots $1$ and $10$ and disability at
slots $2,\ldots,9$, with ten working readings of which eight are positive.
Choose $g_+$ so that their sign confidences agree, with
\[
 0.959375749239844<g_+<0.959375749239846.
\]
Both confidences are $0.876966618\ldots$, whereas their posterior means of the unsigned gain are
$0.678161237\ldots$ and $0.249069192\ldots$.

Each working invocation acquires one body reading and permits the food
protocol of Proposition~\ref{prop:body_memory_value}, or abstention from
that food task. This witness does not acquire or read back the harvest-work
self-label of Definition~\ref{def:joint}: its histories contain body
readings, gate outcomes and the priced clock. The harvest work still
enters the physical ledger. Food switching costs $\kappa/2$ per executed task and
latched recovery costs $\kappa$ per disabled slot. Every working invocation
also pays a declared mandatory processing charge $c_{\rm p}$. A paid clock
signals the slot and horizon; base record, flag, drive and clock supplies
have the same wall-time charges for both latch choices. In this witness,
body-reading erasure and food reset are supplied operations, excluded
from $V(b)$ below rather than assigned zero universal cost.
Proposition~\ref{prop:repair_rates} establishes the corresponding
opposite-repair phenomenon with the terminal reading and reset charges
included. At
$\kappa=2/25$, $c_{\rm p}=0$, the optimal latch differences satisfy
\[
 Q_1(b)-Q_0(b)\in(0.00559146,0.00559147),\qquad
 Q_1(b')-Q_0(b')=-0.004272.
\]
At $\kappa=9/4000$, $c_{\rm p}=3/25$, their differences satisfy
\[
\begin{aligned}
 Q_1(b)-Q_0(b)&\in(0.00587980,0.00587982),\\
 Q_1(b')-Q_0(b')&\in(-0.00432522,-0.00432520).
\end{aligned}
\]
\end{proposition}
\begin{proof}
For a gain $g$ put $a_g=(1+g)/2$ and $b_g=(1-g)/2$. The posterior after
$n$ readings, $k$ positive, is proportional to
$a_g^kb_g^{n-k}$ for sign $+1$ and $b_g^ka_g^{n-k}$ for sign $-1$.
Let $c_{n,k}(x)$ be its positive-sign probability with gains $3/20,x$.
The rational function $c_{2,2}(x)-c_{10,8}(x)$ has negative and positive
values at the displayed rational endpoints. Its denominator is positive,
so continuity gives an exact root. Both histories can be generated by the
same legal policy that holds latch zero until slot ten. Their gate paths
have probabilities $d^2(1-d)^7$ and $(1-d)^9$, independent of the static
instrument; every required reading has positive likelihood. Gate and
clock records therefore leave the count-based posteriors unchanged.

Write $m(b)=\E_b[\sigma g]$ and
\[
 F(x)=\ln2-H_2((1+x)/2),\qquad
 V(b)=\max\{0,F(s|m(b)|)-\kappa/2\}-c_{\rm p},\qquad
 A(b)=\E_b V(b^Y),
\]
where $b^Y$ is the posterior after one further working reading.
Only two future slots remain. If slot eleven works, its optimal
continuation is common to the two current latches. It writes latch zero:
a recovery switched on at slot twelve cannot create a working opportunity
within this horizon. If slot eleven is disabled, the stored latch pays
$\kappa\ell$ there; recovery makes slot twelve working with probability
$r_\ell$, otherwise the same charge is paid once more. Consequently
\begin{equation}\label{eq:short_repair}
 Q_1(b)-Q_0(b)
 =d\bigl[(r_1-d)A(b)-\kappa(2-r_1)\bigr],\qquad
 r_1=d+s(1-2d).
\end{equation}
This is the full Bellman difference for the stated action family.

The numerical intervals follow from rational enclosures, including the
gain-root bracket. For $0\le x<1$ use
\[
 F(x)=\sum_{j=1}^{24}\frac{x^{2j}}{2j(2j-1)}+R_{24}(x),\qquad
 0\le R_{24}(x)\le\frac{x^{50}}{50\cdot49(1-x^2)}.
\]
Rational interval evaluation gives, for the first price pair,
$A(b)\in(0.2236,0.2237)$ and $A(b')=0$.
For the second it gives $A(b)\in(0.1360,0.1361)$ and
$A(b')\in(-0.0954,-0.0953)$; retaining more digits gives the stated intervals.
Substitution in~\eqref{eq:short_repair} proves the strict decisions.
\end{proof}

The histories distinguish how much a future callable task is worth even
when their confidence about its orientation is identical. The processing
charge $3/25$ is a stipulated implementation price, not a universal
erasure lower bound. The first witness needs no such charge. We now use
the close gains and positive drift rates of the continuing instrument,
with reading and food-reset charges from its ledger.

\begin{proposition}[Two moments determine terminal access value]
\label{prop:access_moments}
Immediately before a final working read let $X=\sigma g_G$ have posterior
law $b$, and put $m=\E_bX$ and $v=\E_bX^2$. Use the body reading
$P(Y=y\mid X)=(1+yX)/2$, charge its conditional entropy, and permit either
the posterior-matched food protocol or abstention. For overlap $r$, flip
success $s$ and switching charge $\kappa/2$, the expected working-use
value under these separate conditional-reset prices is
\begin{equation}\label{eq:terminal_access_value}
\begin{gathered}
A(m,v)=-H_2((1+|m|)/2)+
 \sum_{y=\pm1}\frac{1+ym}{2}\,M(m_y),\qquad
m_y=\frac{m+yv}{1+ym},\\
M(z)=\max\{0,S_r(0)-S_r(|z|)
             -H_2((1+s|z|)/2)-\kappa/2\}.
\end{gathered}
\end{equation}
Here $S_r$ is defined in Section~\ref{subsec:J_viability}; a term of zero
observation probability is taken by continuity. A post-record decision
with two wall slots remaining has
\begin{equation}\label{eq:drifting_repair}
 Q_1(b)-Q_0(b)
 =d[(r_1-d)A(m_2,v_2)-\kappa(2-r_1)],\qquad
 r_1=d+s(1-2d),
\end{equation}
where $m_2,v_2$ are the moments of $b\mathsf K^2$.
\end{proposition}
\begin{proof}
Bayes' rule gives $P(Y=y)=(1+ym)/2$ and
$\E[X\mid Y=y]=(m+yv)/(1+ym)$. The food harvest is
$\ln2-H_2((1+s|m_y|)/2)$, its copy-assisted reset costs
$\ln2-S_r(0)+S_r(|m_y|)$, and switching costs $\kappa/2$.
Their difference, optimized against abstention, is $M(m_y)$.
The body-outcome charge is $H_2((1+|m|)/2)$. These terms give $A$.
For the two-slot difference repeat the branch calculation in
\eqref{eq:short_repair}; the instrument undergoes two hidden transitions
before the final reading.
\end{proof}

Two moments suffice for this terminal task, not for the continuing
joint filter. The prices are conditional-entropy limits for specified
separate resets. Base packets, flags, clock, reliable holding and
processor implementation have the common supplied charges of
Table~\ref{tab:body_access_ledger}; the expression is not a total cyclic
apparatus profit. Additional action-dependent costs must be included
in its value before deciding to repair.

\paragraph{One formed model, different decisions.}

The body model carried between invocations is a joint belief over
$(G,\sigma)$, formed by the permitted record history. It has three
uses within this protocol. The food controller acts only when
$M(m_y)>0$; otherwise it withholds permission to spend the switching
and reset resources. The repair controller compares
\eqref{eq:drifting_repair} with zero. A controller that can use a
later invocation for further evidence must also forecast what
today's record will make learnable. These are not three independently
supplied assessments of the observer. They use different functions
of the same acquired belief. A finite horizon and read quota, or a
specified continuing budget, determine which functions are worth
retaining and computing.

\begin{proposition}[Terminal compression need not close enquiry]
\label{prop:enquiry_moments}
Use the same body likelihood as Proposition~\ref{prop:access_moments}
and write $t=\E_bX^3$. Conditional on the present working read,
\begin{equation}\label{eq:enquiry_moments}
 P(Y=y)=\frac{1+ym}{2},\qquad
 m_y=\frac{m+yv}{1+ym},\qquad
 v_y=\frac{v+yt}{1+ym}.
\end{equation}
For a static instrument and two available working invocations,
consider the body-only retained-history protocol: use the terminal
food/abstention rule on each invocation, retain the body record,
and do not use harvest-work labels or the food record as an extra
calibration record. Its two-invocation value, before common
supplied charges, is
\begin{equation}\label{eq:two_call_value}
 A_2(b)=-H_2((1+|m|)/2)
       +\sum_{y=\pm1}\frac{1+ym}{2}
                  [M(m_y)+A(m_y,v_y)].
\end{equation}
Thus $(m,v)$ is sufficient for the last invocation but need not
determine the value of retaining a record for the next. With the
four distinct supported values $\{-g_+,-g_-,g_-,g_+\}$, moments of
orders zero through three determine the entire belief; fewer
need not. Under positive drift, propagate that belief with the
legal event or wall-slot kernel before applying the next likelihood.
\end{proposition}

\begin{proof}
Multiply the prior by $(1+yX)/(1+ym)$ and integrate $X$ and $X^2$.
This gives \eqref{eq:enquiry_moments}. The first invocation pays
its outcome entropy and chooses $M(m_y)$; its retained outcome
determines the prior of the second, whose value is $A(m_y,v_y)$.
Conditioning and averaging prove \eqref{eq:two_call_value}.
For the four-state support, moment recovery is an invertible
Vandermonde system. To see the failure of two-moment closure
without relying on a dimension count, put $a=g_-$, $c=g_+$ and
perturb any strictly positive belief in the stated order by
$e(a/c,-1,1,-a/c)$. Normalization, $m$ and $v$ are unchanged,
whereas $t$ changes by $2ea(a^2-c^2)\ne0$. For example, with
$a=0.90\sqrt{1-0.8^{14}}$, $c=0.98\sqrt{1-0.8^{14}}$,
$s=0.979992$ and $\kappa=0.00225388$, start from
$b=(.02,.08,.10,.80)$ and take $e=.02$.
Both beliefs have $(m,v)\simeq(.765001391,.892279800)$,
but $\sum_yP(y)A(m_y,v_y)$ is respectively
$.181752936$ and $.180874313$. The continuous entropy expressions
in \eqref{eq:terminal_access_value} give this strict difference.
These positive-belief examples establish failure of universal
compression; they are not claimed as the two histories reached
from the uniform prior in Proposition~\ref{prop:repair_rates}.
\end{proof}

The invocation quota in this proposition is real: one body reading
per invocation, no extra calibrated input and no read during a
blackout. Record retention has holding, processing and eventual
reset costs; the gain in \eqref{eq:two_call_value} pays for it
only if it exceeds those prices. An optional additional enquiry
must be compared with foregoing that invocation or using it for
the current task, including its drive, switching and record
charges. The formulas are conditional task values, not a free
probe or a claim of recurrent optimality. In a finite-horizon
control problem the legal belief, remaining quota and available
clock information suffice for the usual Bellman recursion.
If the clock is absent, the disabled duration must instead be
marginalized with \eqref{eq:body_event_kernel}; it cannot be
recovered from unobserved slots. The fixed-support filter and its
precision certificate implement this common model; they do not
supply processor work or a cyclic reset for free.

\begin{proposition}[Paid-repair decisions under positive drift]
\label{prop:repair_rates}
Set $r=0.8^7$, $h=\sqrt{1-r^2}$, $g_-=0.90h$, $g_+=0.98h$,
$\delta=3\times10^{-4}$, $d=1/20$, $s=0.979992$ and
$\kappa=0.00225388$. Start from the uniform four-state prior. Take
a priced clock, horizon $H=12$, and a working decision at slot ten.
Use the terminal protocol of Proposition~\ref{prop:access_moments},
with no harvest-work read-back or further processing surcharge.
At $\eta=2\times10^{-4}$, a lower-confidence history optimally repairs
while a higher-confidence history declines. There is also a strictly
positive $\eta_*$ in
\[
 .00019999976122194292<\eta_*<.00019999976122194293
\]
at which two reachable histories have exactly equal sign confidence
and opposite optimal repairs. The histories and rounded values are
shown in Table~\ref{tab:repair_drift}.
\end{proposition}

\begin{table}[ht]
\centering\small
\begin{tabular}{@{}llrr@{}}
\toprule
Inversion rate & Slot history & Sign confidence & $Q_1-Q_0$\\
\midrule
$0.0002$ & $-\times----++++$ & .969070314 & $+.002191561$\\
$0.0002$ & $++-+------$ & .999973536 & $-.000178194$\\
$\eta_*$ & $-+++\times\times\times+-+$ & .999514604 & $+.000256112$\\
$\eta_*$ & $+--\times-+--+-$ & .999514604 & $-.000845366$\\
\bottomrule
\end{tabular}
\caption{Working body readings are $+$ or $-$; $\times$ is a disabled
slot. All decisions use the same two-slot horizon, prices and action
family. In the last pair confidence is exactly equal at a root in
the stated interval, not merely equal after rounding.}
\label{tab:repair_drift}
\end{table}

\begin{proof}
Hold latch zero through the preceding slots. Each displayed gate path
then has positive probability, independently of the instrument.
All reading likelihoods and instrument transitions are positive.
Starting from the uniform prior, propagate with $\mathsf K$ between
slots and multiply by $(1+y\sigma g_G)/2$ on working readings.
Thus these are reachable beliefs, not assigned posteriors.

The confidence difference of the last two histories is a continuous
function of $\eta$. Outward interval evaluation gives opposite signs
at the two stated rational endpoints, proving existence of an exact
root. Over that whole interval the two differences in
\eqref{eq:drifting_repair} lie respectively in
$(.00025611,.00025612)$ and $(-.00084537,-.00084536)$.
At $\eta=.0002$ the first pair lies in
$(.00219156,.00219157)$ and $(-.00017820,-.00017819)$.
The entropy, square-root and Bayes operations in
\eqref{eq:terminal_access_value} enclose these values directly;
every posterior denominator is strictly positive.
\end{proof}

The enclosures evaluate the displayed Bayes and entropy equations directly,
rather than using sampled trajectories to decide the four histories.

The first pair needs no tuned parameter and rules out an increasing
confidence-threshold repair rule at the continuing instrument's own
rates. The second rules out sign confidence as a sufficient decision
statistic in a nearby positive-rate model. It does not assert a unique
root. The restoring reader assigns more weight to the higher gain;
the declining reader expects less value from a working invocation.
Confidence in orientation and the value of accessible tasks are
different properties of the same acquired history.

The finite horizon is material. A continuing observer must also price
evidence gained on later invocations, lost evidence during disability,
and the registers that carry its model. Recurrent optimal maintenance
without an elapsed-time counter remains a stronger target. The two
moments above explain a terminal decision without adding history
as an independent reward.

\section{Consequences for the other links}
\label{subsec:J_consequences}

\emph{Paper~II.} Let the agent add a known push $p_t$, independent of the input, through an actuator with gain $k$, so that $y_t=\rho\,u\cdot x_t+k\rho\,u_A\cdot p_t+\xi_t$, with one in-frame read per step. If sensing and acting share one coupling ($u_A=u$), the regressions of $y$ on the push and on the in-frame read, with the residual variance, determine $u$ including its background, $\rho$, $k$ and $s$; the blind spot of Paper~II disappears with a single read per step when $u\cdot v\ne0$, $k\rho\ne0$ and sign conventions for $k,\rho$ are fixed, and the remaining symmetry between actuator gain and input scale is fixed by the known input law. If the actuator coupling is independent of the sensing coupling, pushes determine $k\rho u_A$ and nothing about the background of $u$. In a simulation with $D=6$ the shared coupling recovers $\rho$ and $k$ within $2\%$ and the direction of $u$ within $0.016$; with an independent coupling the direction error is $1.3$. A shared sensing/action frame supplies this constructive route. At $u\perp v$ the scale parameters remain confounded; varying the read direction repairs the degeneracy.

\emph{HAFF Paper~C.} The viability bound of an agent depends on it only through the product of its sensing and control influences~\cite{Liu2026HAFF_C}, which is invariant under the sign symmetry, while viability itself is not: an inverted agent has the same influences and fails with probability $\tfrac12(1+(1-2\nu)(1-2\mu)(1-2q))$. Positive influences can be certified from the agent's own records; viability cannot be certified without a standard.

\emph{Several observers.} Applying Theorem~\ref{thm:gauge} with the sign symmetry acting on two agents at once, their relative sign is invariant and detectable from agreement, and a common inversion is not. The common error that agreement can never expose, which Foundation~II measures~\cite{Liu2026FII}, is the direction shared by the observers' standards.


% ============================================================
\section{Discussion}
\label{sec:discussion}

\subsection{The memory--self--introspection route}
\label{subsec:logic_chain}



The shared channel, loss metric and frame residual were specified in Section~\ref{subsec:sequence}. The body adds a further use of the learned model: deciding which future access is worth restoring. Proposition~\ref{prop:repair_rates} proves that sign confidence can select the wrong repair even at positive drift rates; Propositions~\ref{prop:access_moments}--\ref{prop:enquiry_moments} identify the extra retained information needed for terminal and prospective tasks. The body--permission coupling carries a relational standard. Wiring, transition laws and resource supplies remain specified inputs, and joint relabelling preserves the absolute-orientation obstruction of Theorem~\ref{thm:gauge}. A cyclic processor with holding and preparation costs, and a continuous self-frame with matrix-valued loop gains, remain stronger targets.

\subsection{What is established and what is open}
\label{subsec:claims}

Within the stated models the results are proved, and several are sharp. The detection bound holds for every test; it is reached within a factor of about $1.3$ when the drift path is known, and its order is reached by a simple residual-variance test under Brownian drift. The scalar quadratic drift law holds for every schedule and estimator; Gaussian prior/readings and controllable real precision attain its optimum, so the price of keeping up is known exactly to leading order, with its constant. The task-risk formula of Section~\ref{sec:loop_world} is exact for every joint law of the reports.

The main open problems are these. The drift law is proved for a scalar parameter and a Gaussian random walk; vector parameters, non-Gaussian increments (for which the same proof gives a weaker bound with the increment variance replaced by the inverse Fisher information of the increment law) and smoother drift (Conjecture~\ref{conj:smooth}) remain. Paper~II conjectures that probe calibration of a $D$-channel frame is rate optimal for spherical Brownian/local isotropic diffusion, not for every rotation-invariant drift law~\cite{Liu2026TDOME_II}; the telescoping of Theorem~\ref{thm:driftlaw} applied per background direction is a possible route, but the directions of a frame on a sphere are not independent random walks. No rule is given that learns the task risk of Proposition~\ref{prop:task_risk} online, and the gain proxy of Section~\ref{sec:loop_world} is not shown to be the best measure of reliability. Fisher information is the natural local quantity for these questions; nothing here shows that a physical system must compute it.

The positive-rate repair witness establishes a finite-horizon mechanism,
not an optimal recurrent policy. Its extension must combine continuing
maintenance, a finite processor and the cost of acquiring its decision
information. A physical cycle must supply holding, processing, reference
preparation and heat transport together. BRAIN's component-maintenance
model~\cite{Liu2026BRAIN} treats the corresponding question when learned
relations span several ports; it has a different repair clock and
cannot inherit the single-body event kernel without an explicit mapping.

\subsection{Function and experience}
\label{subsec:function}

Reflexivity in this paper means second-order control: a system that monitors and adjusts its own monitoring. The results say what such control can detect, what it costs and what it buys in a loop. The access gate also makes the reliability of future self-knowledge depend on physical maintenance. Its working and disabled states are operational permissions, not claims about sleep, death or experience. They do not show that a calibrated self-model is, or produces, experience. The Consciousness series treats the relation between functional self-models and first-person organization with its own arguments~\cite{Liu2026CIII}.

% ============================================================
\appendix

\section{Worked example}
\label{sec:example}

\subsection{Detecting drift in Paper~II's two-channel reader}
\label{subsec:example_detect}

Take Paper~II's reader with $D=2$, $\rho=1$, $s=0.5$ and one record per step ($r=1$). Then $\kappa=36$ for in-frame records and $F_\perp(0)=0.8$ for probes. For errors $\alpha=\beta=0.05$, $d(0.95\Vert0.05)=2.650$. Table~\ref{tab:detect} gives the smallest horizon at which the exact relative entropies of Corollary~\ref{cor:three_records}, summed over the records, reach $d$. By Proposition~\ref{prop:detect} no test with these errors can succeed earlier.

\begin{table}[t]
\centering
\small
\begin{tabular}{lccc}
\toprule
Linear drift $\varepsilon$ & $10^{-2}$ & $3\times10^{-3}$ & $10^{-3}$\\
\midrule
In-frame records, exact (asymptotic) & 54 (49) & 137 (129) & 324 (312)\\
Probe records, exact (asymptotic) & 58 (58) & 130 (130) & 271 (271)\\
\midrule
Brownian drift $\sigma$ & $10^{-2}$ & $10^{-3}$ & $10^{-4}$\\
\midrule
In-frame records, exact (asymptotic) & 17 (11) & 68 (53) & 282 (245)\\
Residual-variance test reaches errors $0.05$ & 84 & 356 & 1565\\
Frozen self-model fails at $t^*=\ln2/\sigma$ & 69 & 693 & 6931\\
\bottomrule
\end{tabular}
\caption{Detection times in steps for Paper~II's reader ($\rho=1$, $s=0.5$, one record per step, errors $0.05$). The rows labelled ``exact (asymptotic)'' are necessary times: ``exact'' sums the exact relative entropies, with the unwrapped angle normal of variance $2\sigma t$ under Brownian drift, and ``asymptotic'' is the formula of Proposition~\ref{prop:detect}. The test row is the smallest horizon at which the residual-variance test of Corollary~\ref{cor:before_failure}, with the exact $\chi^2$ threshold, reaches both errors $0.05$ (Monte Carlo).}
\label{tab:detect}
\end{table}

For known linear drift, the likelihood-ratio test reaches errors $0.05$ at about $1.3$ times the necessary horizon (4000 repetitions per point). The path-free residual-variance test needs roughly five times the Brownian lower bound; it precedes frozen-model failure at $\sigma=10^{-3},10^{-4}$, but not at $10^{-2}$. Its $T\sigma^{2/3}$ approaches the predicted constant $3.28$. Near calibration the in-frame natural gradient is $0.0556/\theta$, whereas a probe fraction $p$ supplies metric $0.8p$ at calibration.

\subsection{The local Gaussian comparison for a drifting gain}
\label{subsec:example_price}

Let the gain of the gain--bias channel drift with $g\approx0.55$, $b=0$, so $J=G=1.4337$, and let references cost $w=k_BT\ln2$. Table~\ref{tab:price} gives the local Gaussian comparison of Remark~\ref{cor:thermo}, its leading term $C$, the optimal reference rate and the optimal memory weight.

\begin{table}[t]
\centering
\small
\begin{tabular}{lcccc}
\toprule
$\nu$ per use & Floor ($k_BT$) & $C$ ($k_BT$) & $\rho^*$ (references per use) & $\alpha^*$\\
\midrule
$10^{-2}$ & 0.05511 & 0.05515 & 0.0265 & $1.95\times10^{-3}$\\
$3\times10^{-3}$ & 0.02471 & 0.02471 & 0.0119 & $3.9\times10^{-4}$\\
$10^{-3}$ & 0.01188 & 0.01188 & 0.0057 & $9.1\times10^{-5}$\\
$10^{-4}$ & 0.002560 & 0.002560 & 0.0012 & $4.2\times10^{-6}$\\
\bottomrule
\end{tabular}
\caption{The local Gaussian comparison for a drifting gain ($g=0.55$, $b=0$, $w=k_BT\ln2$). ``Floor'' is the Gaussian quadratic comparison $\min_Ph(P)$ of Theorem~\ref{thm:driftlaw} for the quadratic loss; $C$ is the leading term of Corollary~\ref{cor:driftlaw}. For Paper~I's divergence loss the entries are leading-order values, accurate for $\nu\lesssim10^{-3}$. A record read through the same channel is worth $0.160\,k_BT$ (Paper~I).}
\label{tab:price}
\end{table}

At $w=J=c=1$ and $\nu=10^{-2}$, average-cost value iteration on a variance grid agrees with the Gaussian floor to relative error $1.3\times10^{-5}$; 4000 Kalman trajectories agree within Monte Carlo error. Integer binomial reads provide a distinct comparison: at $w=10^{-4}k_BT$, $\nu=3\times10^{-3}$, the best tested grid policy costs $1.10$ times the relaxed floor. For Corollary~\ref{cor:c3}, $q=10^{-4}$, $r=v=1$, $w=10^{-3}$ gives optimum $0.0054878$. The source map identifies the full scans and fixed-seed outputs.


\section{A drifting instrument in a closed loop}\label{sec:loop_world}

Two finite bank models couple the estimate $(\hat g,\hat b)$ to an action that changes the next world reading. Their instruments, baths and results remain separate. They specify their own assumptions and do not use (C1)--(C3).

\subsection{The world and the observation map}\label{subsec:bank}

The bank has $M$ binary variables $z_j\in\{\pm1\}$ with mean $x=M^{-1}\sum_jz_j$. In each epoch every variable is updated independently: it keeps its value with probability $a$, is replaced by a variable biased toward the agent's command $u$ with probability $c$, and is replaced by a variable biased toward the environmental input $d$ with probability $r$, where $a+c+r=1$ and the biased variable is $+1$ with probability $(1+v)/2$ for the bias $v=u$ or $v=d$. For $|u|,|d|\le1$ this is a convex combination of the identity and two resets, hence a legal channel (Foundation~III, section on legal finite feedback~\cite{Liu2026FIII}). The probability of $+1$ is therefore $(1+az_j+cu+rd)/2$, and
\begin{equation}\label{eq:bank_step}
x_{k+1}=a\,x_k+c\,u_k+r\,d_k+\zeta_k,\qquad \mathbb E[\zeta_k\mid\text{past}]=0,\quad \mathbb E[\zeta_k^2\mid\text{past}]\le M^{-1}.
\end{equation}
The count of $+1$ variables is updated exactly from two binomial draws. The simulations hold that count for efficiency and never give it to the agent. We use $(a,c,r)=(0.65,0.25,0.10)$.

An instrument with gain $g$ and bias $b$, $|g|+|b|\le1$, reads a uniformly random variable and returns $y=\pm1$ with probability $(1+gz_j+b)/2$. The index is not returned. A reference read presents an input known to be $+1$ or $-1$ and returns $+1$ with probability $(1\pm g+b)/2$. Gain and bias are fixed within an epoch and drift between epochs. From $n_s$ world reads with mean $\bar y$ and the current estimates the agent forms
\begin{equation}\label{eq:xhat_loop}
\hat x=Q\!\left(\frac{\bar y-\hat b}{\max(\hat g,\,0.30)}\right),
\end{equation}
where $Q$ clips to $[-1,1]$ and rounds to a grid of $1025$ points, and it acts with $u=-h\hat x$, $h=0.8$. The total number of reads per agent and epoch is $512$. World loss is $x^2/4$, averaged over trajectories and over epochs after a burn-in of $200$.

\subsection{What reference reads measure}\label{subsec:reference_fisher}

\begin{proposition}[Reference reads]\label{prop:reference_reads}
Let $p\sim\mathrm{Bin}(n,\pi_+)$ and $m\sim\mathrm{Bin}(n,\pi_-)$ be independent, with $\pi_\pm=(1\pm g+b)/2$ and $|g|+|b|<1$. Then $(\hat g,\hat b)=\bigl((p-m)/n,\;(p+m)/n-1\bigr)$ is unbiased with covariance
\[
\Sigma_{\mathrm{ref}}=\frac1{2n}\begin{pmatrix}1-g^2-b^2&-2gb\\-2gb&1-g^2-b^2\end{pmatrix},
\]
which is the inverse of the Fisher information matrix of the $2n$ reads about $(g,b)$. The estimator is therefore efficient.
\end{proposition}
\begin{proof}
$\mathrm{Var}\,p=n(1-(g+b)^2)/4$ and $\mathrm{Var}\,m=n(1-(b-g)^2)/4$, which gives the covariance by linearity. The pair $(p/n,m/n)$ is the maximum-likelihood estimator of the mean parameters $(\pi_+,\pi_-)$ of a product of binomials and attains the Cram\'er--Rao bound $\mathrm{diag}(\pi_\pm(1-\pi_\pm))/n$. The map to $(g,b)$ is linear and invertible, and efficiency is preserved under it.
\end{proof}

The frequency estimator above is unbiased and efficient for mean-square
error, but its boundary values make the plug-in expected divergence infinite
at every finite sample size. To use Paper~I's calibrated-work expansion,
replace each estimated response probability $k/n$ by
$(k+1/2)/(n+1)$, or use another regularization satisfying its tail condition.
The leading covariance is unchanged, and the $1/n$ work loss is
then an asymptotic statement for that regularized estimator.

For that regularized estimator, Paper~I computes the leading calibrated loss: acting on the estimated readout model costs $k_BT(1+g^2/(1-b^2))/(4n)$ per use to leading order, $0.00254\,k_BT$ for $n=128$ at $g=0.55$, $b=0$. A reference read costs a world read, so $n_r$ is a budget. The estimates are also kept in memory. A persistent memory with weight $\alpha$ on the new reference estimate, $\hat\theta_k=Q[(1-\alpha)\hat\theta_{k-1}+\alpha\tilde\theta_k]$ for $\theta\in\{g,b\}$, trades noise against lag.

\begin{proposition}[Memory under drift]\label{prop:memory_drift}
Let $0<\alpha\le1$, $\beta=1-\alpha$ and $e_k=\hat\theta_k-\theta_k$. Suppose the reference estimate has the form $\tilde\theta_k=\theta_k+\xi_k$, where $\theta_k$ is determined by the information available before the reference reads of epoch $k$ (a deterministic drift qualifies, and the true instrument state may enter without giving the observer access to it) and $\xi_k$ has conditional mean zero given that information and variance at most $\sigma^2$, that $\theta_k\in[-1,1]$ and $|\theta_k-\theta_{k-1}|\le\nu$, and that $Q$ clips to $[-1,1]$ and rounds to a grid of spacing $\Delta$. Then
\[
\|e_k\|_{L^2}\le\sqrt{\bigl(\beta(\|e_{k-1}\|_{L^2}+\nu)\bigr)^2+\alpha^2\sigma^2}+\frac\Delta2,
\qquad
\|e_k\|_{L^2}\le X_\sharp+\beta^k\bigl(\|e_0\|_{L^2}-X_\sharp\bigr)_+,
\]
with $X_\sharp=\dfrac{\Delta/2+\beta^2\nu+\sqrt{\beta^2(\Delta/2+\nu)^2+(1-\beta^2)\alpha^2\sigma^2}}{1-\beta^2}$. If $0<\alpha<1$, the noise variance is exactly $\sigma^2$, and the drift is the idealised linear drift $\theta_k=\theta_0+\nu k$ without rounding or clipping, the steady-state error has mean $-\nu(1-\alpha)/\alpha$ and variance $\alpha\sigma^2/(2-\alpha)$. Hence the persistent memory has smaller mean-square error than a memoryless estimate (weight $\alpha=1$, error variance $\sigma^2$) exactly when
\[
\nu^2<\frac{2\alpha^2\sigma^2}{(1-\alpha)(2-\alpha)}.
\]
\end{proposition}
\begin{proof}
Put $y_k=\beta\hat\theta_{k-1}+\alpha\tilde\theta_k$, so that $y_k-\theta_k=\beta e_{k-1}-\beta(\theta_k-\theta_{k-1})+\alpha\xi_k$. The noise $\xi_k$ is orthogonal in $L^2$ to the first two terms, which are determined before it, so $\|y_k-\theta_k\|^2\le\beta^2(\|e_{k-1}\|+\nu)^2+\alpha^2\sigma^2$. Clipping to an interval that contains $\theta_k$ does not increase the distance to $\theta_k$, and rounding adds at most $\Delta/2$; this is the first inequality. Write $\varphi(X)$ for its right-hand side, which is increasing with $\varphi'(X)=\beta^2(X+\nu)/\sqrt{\beta^2(X+\nu)^2+\alpha^2\sigma^2}\le\beta$. The larger root $X_\sharp$ of $(X-\Delta/2)^2=\beta^2(X+\nu)^2+\alpha^2\sigma^2$ is the displayed expression, and it satisfies $X_\sharp\ge\Delta/2$, so $\varphi(X_\sharp)=X_\sharp$ and $\varphi(X_\sharp+t)\le X_\sharp+\beta t$ for $t\ge0$. The second inequality follows by induction. For linear drift without rounding or clipping the recursion is linear, $e_k=\beta e_{k-1}-\beta\nu+\alpha\xi_k$, which gives the stated mean and variance. The mean-square comparison is $\alpha\sigma^2/(2-\alpha)+\nu^2(1-\alpha)^2/\alpha^2<\sigma^2$.
\end{proof}

For Model~A's $\alpha=0.25$ and $n_r=128$ reads per sign, the unrounded linear-drift comparison gives $\nu^*\approx0.0161$ per epoch. It places the crossover at the scale of the observed four-to-eight-fold drift bracket; it omits the loop, gain floor and periodic drift. The Gaussian-transition optimum of Theorem~\ref{thm:driftlaw} instead uses the stationary Kalman gain and real precision. A constant-weight estimator under linear drift has the different $\nu^{2/5}$ price of Remark~\ref{rem:driftlaw_scope}.

The world is read through the same estimates. The next statement bounds the error of the decoded state in terms of the errors of the estimates.

\begin{proposition}[Error of the decoded state]\label{prop:decode_error}
Let $g\ge g_{\min}>0$ and let $\bar y$ be the mean of $n_s$ readings of the bank, which are independent given the state $x$. Let $\hat x=Q_x\bigl[\operatorname{clip}\bigl((\bar y-\hat b)/\max(\hat g,g_{\min})\bigr)\bigr]$, where $Q_x$ rounds to a grid of spacing $\Delta_x$ that contains $\pm1$. Then
\[
\|\hat x-x\|_{L^2}\le\min\!\left(2,\ \frac{n_s^{-1/2}+\|\hat g-g\|_{L^2}+\|\hat b-b\|_{L^2}}{g_{\min}}+\frac{\Delta_x}2\right).
\]
\end{proposition}
\begin{proof}
Write $\bar y=gx+b+\nu$ with $\mathbb E[\nu\mid x]=0$ and $\mathbb E[\nu^2\mid x]\le n_s^{-1}$, and $\hat g'=\max(\hat g,g_{\min})\ge g_{\min}$. Then $(\bar y-\hat b)/\hat g'-x=[(g-\hat g')x+(b-\hat b)+\nu]/\hat g'$. Since $g\ge g_{\min}$ one has $|g-\hat g'|\le|g-\hat g|$, and $|x|\le1$, so the numerator has $L^2$ norm at most $\|g-\hat g\|+\|b-\hat b\|+n_s^{-1/2}$, and the division by $\hat g'\ge g_{\min}$ gives the bound before clipping. Clipping to $[-1,1]\ni x$ does not increase the distance to $x$, and rounding adds at most $\Delta_x/2$. Both $\hat x$ and $x$ lie in $[-1,1]$, so the error is at most $2$.
\end{proof}

Model~A has $\nu_g\le0.005357$, $\nu_b\le0.006585$ and $g\ge0.30$. With $\alpha=0.25$, $n_r=128$, $\sigma^2=1/(2n_r)$ and grid spacing $2/1024$, Proposition~\ref{prop:memory_drift} gives limiting $L^2$ errors at most $0.0351$ and $0.0376$. Proposition~\ref{prop:decode_error}, with $n_s=256$, then bounds decoded-state error by $0.4518$, compared with observed epoch-averaged error $0.135$. These are expectation bounds for the stated drift, not simultaneous guarantees or bounds for the zero-gain fault.

\subsection{One observer: the loop closes on the world}\label{subsec:model_A}

\paragraph{Model A.} A single agent, a drifting instrument and a bath $d_k=0.60+0.25\sin(2\pi k/177)$. Normal drift is
\begin{align*}
g(k)&=0.55+0.17\sin(2\pi k/420)+0.03\sin(2\pi k/67),\\
b(k)&=0.14\sin(2\pi k/300)+0.025\sin(2\pi k/43),
\end{align*}
and a drift factor $s$ replaces $k$ by $sk$. The static instrument has $g=0.55$, $b=0$ and the same bath. Three policies share the interface and the budget of 512 reads. \emph{Frozen}: 512 world reads, calibrated once from $2048$ reads per sign at the start and never again. \emph{Memoryless}: 256 world reads and 128 reads per sign each epoch, using only that epoch's references ($\alpha=1$). \emph{Persistent}: the same reads, $\alpha=0.25$. Each policy pays the 4096 start-up reads. The agent does not know the drift functions.

Table~\ref{tab:loop} retains the static, normal-drift and fast-drift comparisons. A separate implementation and larger normal-drift run agreed within $1.7$ combined standard errors; their full scans are identified in the source map.

\begin{table}[htbp]
\centering
\begin{tabular}{lccc}
\hline
Instrument & Frozen & Memoryless & Persistent \\
\hline
static instrument & 3.426$\pm$0.036 & 3.589$\pm$0.009 & 3.520$\pm$0.008 \\
drift $\times$1 & 5.299$\pm$0.043 & 3.652$\pm$0.010 & 3.591$\pm$0.010 \\
drift $\times$4 & 4.642$\pm$0.040 & 3.654$\pm$0.010 & 3.622$\pm$0.009 \\
drift $\times$8 & 4.543$\pm$0.042 & 3.643$\pm$0.008 & 3.802$\pm$0.008 \\
\hline
\end{tabular}
\caption{World loss $x^2/4$ ($\times10^{-3}$, $\pm$ one standard error over 160 independent trajectories, 1200 epochs, burn-in 200).}
\label{tab:loop}
\end{table}

At normal drift, persistent calibration lowers loss by $1.7\%$ relative to memoryless calibration and $32.2\%$ relative to the frozen policy. Freezing is best for the static instrument. The persistent-minus-memoryless difference changes from $-0.032\pm0.013$ at four-fold drift to $+0.159\pm0.012$ at eight-fold drift ($\times10^{-3}$). Memory therefore helps only within a drift range at this weight and read allocation.

\subsection{Two observers: reliability and authority}\label{subsec:model_B}

\paragraph{Model B.} Two agents read the same bank through different instruments (the instruments and the public event of Section~5 of Foundation~II~\cite{Liu2026FII}, here with $d_k=0.6\,c_k$), each with the calibration of Section~\ref{subsec:reference_fisher} and the same budget of 512 reads. Each agent forms $\hat x_i$ by~\eqref{eq:xhat_loop} and issues $u_i=-h\hat x_i$. The world receives the joint command $u=wu_1+(1-w)u_2$, so $|u|\le h$ whenever $w\in[0,1]$. The agents may also move their bias estimates toward the partner's report, either with equal weight or weighted by $\mathrm{clip}(\hat g_j/0.4,0,1)^2$, which requires exchanging $\hat g$ as well as $\hat x$ (22 bits per sender and epoch instead of 11, so 44 for the pair instead of 22). The fixed rule is $w=\tfrac12$. The convex rule is $w=s_1/(s_1+s_2)$ with $s_i=\mathrm{clip}(\hat g_i/0.4,0,1)^2$, and $w=\tfrac12$ if both vanish. In the convex rule the more reliable agent receives more of the authority, and the total never exceeds the cap $h$.

\paragraph{One-step task risk.} Let the reports be $\hat{\boldsymbol x}=x\mathbf 1+\epsilon$ with errors $\epsilon$ of any joint law, and let the weights $w\ge0$, $\mathbf 1^{\!\top}w=1$, be determined by the information $\mathcal F$ available when they are chosen, which may include the reference reads of the epoch. Write $q=ch$, $z=(a-q)x+rd$, $R=\mathbb E[\epsilon\epsilon^\top\mid\mathcal F]$ (an uncentred second moment) and $b=\mathbb E[z\epsilon\mid\mathcal F]$. Weights are compared with the conditional laws of the reports and of the world held fixed.

\begin{proposition}[Task risk of a weight vector]\label{prop:task_risk}
Suppose $|u|\le1$ and $|d|\le1$, so that all branch probabilities of the bank step lie in $[0,1]$, and the action is the linear rule $u=-h\,w^\top\hat{\boldsymbol x}$. Then
\[
\mathbb E\!\left[\tfrac14x'^2\,\middle|\,\mathcal F\right]
=K_{\mathcal F}+\frac{1-M^{-1}}4\left(q^2\,w^\top Rw-2q\,b^\top w\right),
\]
where $K_{\mathcal F}=\tfrac14\bigl[(1-M^{-1})\mathbb E[z^2\mid\mathcal F]+M^{-1}\mathbb E[1-a^2+a^2x^2\mid\mathcal F]\bigr]$ does not depend on $w$.
\end{proposition}
\begin{proof}
Given $x$, $u$ and $d$ the variables move independently, so $\mathbb E[x'\mid x,u,d]=ax+v$ and $\mathrm{Var}(x'\mid x,u,d)=M^{-1}(1-a^2-v^2-2axv)$ with $v=cu+rd$. Hence $\mathbb E[x'^2\mid x,u,d]=(1-M^{-1})(ax+v)^2+M^{-1}(1-a^2+a^2x^2)$. With $u=-h(x+w^\top\epsilon)$ one has $ax+v=z-q\,w^\top\epsilon$. Taking the conditional expectation over $\epsilon$ gives the formula.
\end{proof}

The proposition needs no assumption that the errors are unbiased, centred, or independent between the two agents. Quantisation and the gain floor can sit inside $\epsilon$. If the action is saturated or non-linear the formula no longer applies, and for $q=0$ the weights do not change the one-step risk. When the errors are unbiased given $x$, $d$ and $\mathcal F$, $b=0$ and $R$ is the error covariance. The risk is a convex quadratic in $w$. For two agents with $w=(t,1-t)$, $q>0$ and $D=R_{11}+R_{22}-2R_{12}>0$, its minimiser over $t\in[0,1]$ is
\begin{equation}\label{eq:t_star}
t_*=\mathrm{clip}\!\left(\frac{R_{22}-R_{12}+(b_1-b_2)/q}{D},\,0,\,1\right).
\end{equation}
If $R=\mathrm{diag}(v_1,v_2)+\tau^2\mathbf 1\mathbf 1^\top$ and $b=0$, the best combined error is $\tau^2+v_1v_2/(v_1+v_2)$: weighting removes the private part of the error and cannot touch the common part.

\begin{proposition}[Estimation-optimal and task-optimal weights differ]\label{prop:two_optima}
Take $M=20$, $x=0.2$ (twelve variables at $+1$), $(a,c,r)=(0.65,0.25,0.10)$, $h=0.8$, $d=-0.7$. Let $\epsilon_1$ be $0$ or $0.2$ and $\epsilon_2$ be $-0.1$ or $0.1$, each with probability $\tfrac12$ and independent. The weight that minimises the joint estimation error is $t=\tfrac13$; the weight that minimises the next-step world loss is $t=\tfrac23$.
\end{proposition}
\begin{proof}
The estimation error is $t^2\,0.02+(1-t)^2\,0.01$, minimised at $t=\tfrac13$. Here $q=0.2$, $z=0.02$, $R=\mathrm{diag}(0.02,0.01)$ and $b=(0.002,0)$, so~\eqref{eq:t_star} gives $t_*=\tfrac23$. In every branch $|u|\le0.32$, so the step is legal. The values in Table~\ref{tab:two_optima} come from enumerating all four error branches and both binomial draws exactly, and they agree with Proposition~\ref{prop:task_risk}.
\end{proof}

\begin{table}[htbp]
\centering
\begin{tabular}{lccc}
\hline
Allocation rule & $t$ & Joint estimation error & Next-step world loss \\
\hline
Smallest joint estimation error & 0.3333 & 0.006667 & 0.007525 \\
Inverse centred variance & 0.5000 & 0.007500 & 0.007501 \\
Smallest next-step world loss & 0.6667 & 0.010000 & 0.007493 \\
\hline
\end{tabular}
\caption{The example of Proposition~\ref{prop:two_optima}, by exact enumeration.}
\label{tab:two_optima}
\end{table}

Agent~1's mean error $0.1$ offsets $z=0.02$, making more authority useful despite worse estimation. Task-optimal allocation needs the effects of action as well as sensor reliability. The observable gain proxy below is not the covariance $R$; no online learner for the true-error quantities $R,b$ is supplied.

The linear term $-2q\,b^{\top}w$ also marks where Paper~I's metric applies. Paper~I prices a calibration error by $\tfrac12k_BT\,\delta^{\top}G\delta$ when the downstream protocol is matched to the posterior the agent assumes, and then a calibrated self-model is a stationary point of the loss. With a fixed linear action rule it need not be one, and a biased report can lower the world loss, as in Proposition~\ref{prop:two_optima}.

\paragraph{Four cells.} We cross the two interventions that can be applied separately: the update rule (equal or weighted alignment) and the authority rule (fixed or convex). One agent's instrument is damaged for 100 epochs at a random start: a bias fault raises its bias by a nominal $0.30$ before the legal clip, a gain fault sets its gain to $0.05$ or $0$ (``zero gain'': the current reading carries no information about $x$). Table~\ref{tab:cells} gives the increase in world loss during the window relative to a matched run without the fault. The commands in all runs are $u_i=-h\hat x_i$ with $\hat x_i$ on the 1025-point grid, so $|u_i|\le h$ exactly; the largest $|u|$ applied over 105 runs was 0.8000 and no epoch exceeded the cap. Every cell is charged the same message budget, 22 bits per sender and epoch ($\hat x$ and $\hat g$) and 44 for the pair, although the unweighted, fixed cell does not use $\hat g$.

\begin{table}[htbp]
\centering
\begin{tabular}{lccc}
\hline
Policy & Bias fault & Gain $0.05$ & Gain $0$ \\
\hline
Unweighted, fixed $w=\tfrac12$ & +0.01$\pm$0.01 & +2.05$\pm$0.02 & +2.53$\pm$0.02 \\
Unweighted, convex authority & +0.03$\pm$0.01 & +1.11$\pm$0.02 & +1.30$\pm$0.02 \\
Weighted, fixed $w=\tfrac12$ & +0.01$\pm$0.01 & +1.32$\pm$0.02 & +1.58$\pm$0.02 \\
Weighted, convex authority & +0.01$\pm$0.01 & +0.34$\pm$0.02 & +0.31$\pm$0.02 \\
References only, no coupling & +0.00$\pm$0.01 & +1.46$\pm$0.02 & +1.80$\pm$0.02 \\
\hline
\end{tabular}
\caption{Increase of world loss in the fault window ($\times10^{-3}$, mean of three seed blocks of 800 trajectories $\pm$ combined standard error; the different policies use different random draws, so the comparison between rows is not a paired test).}
\label{tab:cells}
\end{table}

Weighted alignment limits a failed report's influence on the other estimate; convex authority limits its influence on the world while preserving the command cap. Together they reduce the zero-gain fault increase by $87.8\%$ relative to equal alignment and fixed authority. The interaction, $-0.042\pm0.035$ ($\times10^{-3}$), is unresolved. Fault-free losses range from $3.156$ to $3.167$ ($\times10^{-3}$), without a claim of equality.

\subsection{What the two models give}\label{subsec:loop_scope}

Both models are finite banks with a declared instrument family, a declared bath and a declared cost vector (reads, start-up references, bits of memory and message). Model~A shows a self-model built from reference reads improving the agent's effect on its world, with a memory whose useful length is limited by the drift. Model~B shows that reliability can be used in two separate places, in what an agent believes about its partner and in how much of the world it is allowed to move, and that the second place accounts for the larger part of the reduction when a sensor fails. The two models use different instruments and baths and are not one experiment.


% ============================================================


\section{Simulation protocols for Section~\ref{sec:loop_world}}\label{app:loop_protocol}

\paragraph{Model A.} Use the parameters and policies of Section~\ref{subsec:model_A} and start with $n_+=\lfloor0.6M\rfloor$, $M=4096$. Each epoch reads the pre-action bank and references, updates the estimates, computes~\eqref{eq:xhat_loop}, then applies~\eqref{eq:bank_step}. The agent receives only read counts and its stored estimates. It retains 22 estimate bits and an 11-bit command index, pays 4096 startup references, 256 references per adaptive epoch and 4096 bank-step applications. Counters, index generation and arithmetic workspace are additional resources; no heat or work is inferred from these counts. Run lengths and errors are specified in Table~\ref{tab:loop}.

\paragraph{Model B.} The same bank and step, with the public event $c_k$ (switching with probability $0.01$ per epoch) and $d_k=0.6\,c_k$. The instruments are
\begin{align*}
g_i(k)&=0.55+0.17\sin(2\pi k/420+\phi_i)+0.03\sin(2\pi k/67+\phi_i),\\
b_i(k)&=0.10\sin(2\pi k/250)+0.03\sin(2\pi k/53)\\
&\quad+0.10\sin(2\pi k/300+\psi_i)+0.025\sin(2\pi k/43+\psi_i),
\end{align*}
with $\phi=(0,1.7)$ and $\psi=(0.4,2.9)$, clipped to $|b|\le0.98-|g|$. \begin{sloppypar}Each agent makes $512$ reads per epoch, of which $2n_r$ with $n_r=16$ are reference reads, and starts from 2048 reference reads per sign. Order within an epoch: both agents read the world and references on the same pre-action state; each updates its own calibration and forms $\hat x_i$; one synchronous exchange; the commands $u_i=-0.8\,\hat x_i$ and the joint command are formed from the values before the exchange; the world takes one joint step; the alignment update for the next epoch uses the exchanged values. Equal alignment: $\hat b_i\leftarrow Q(\hat b_i+0.25\,e_i)$ with $e_i=\bar y_i-(\hat g_i\hat x_j+\hat b_i)$. Weighted alignment multiplies the step by $\mathrm{clip}(\hat g_j/0.4,0,1)^2$. The fault window lasts 100 epochs from a start drawn uniformly from $[800,2400)$ for each trajectory; faulted and no-fault runs share it. Runs of 3000 epochs, burn-in 200, 800 trajectories per run, three seed blocks. The world and each agent draw from separate random streams addressed by seed, epoch and event, so a fault in one instrument does not shift the draws of the world or of the other agent; because the agents act on the world, faulted and no-fault worlds differ afterwards.\end{sloppypar}

\paragraph{Reproducibility.}
The supplemental source map lists checks by manuscript table, proposition
or equation label, with named scripts, archived outputs and hashes,
including the two Model~A implementations and Model~B.
Checks cover detection, drift prices, gauge symmetry, sign inversion,
body moments, finite-block margins, continuing filtering and repair.
Their individual assumptions and seeds are retained with the scripts;
one experiment's output is not a certificate for a different model.

% ============================================================
\section{Thermodynamic implementation inputs}
\subsection{Thermodynamic length}
\label{subsec:thermo_length}

Let $\lambda(t)$, $t\in[0,\tau]$, be a protocol in control-parameter space. For conjugate forces $F_i=-\partial_i H$, define the friction tensor by
\[
\zeta_{ij}(\lambda)=\beta\int_0^\infty
  \langle\delta F_j(0)\delta F_i(t)\rangle_{\lambda}\,dt.
\]
The thermodynamic length is $\mathcal L=\int_0^\tau\sqrt{\zeta_{ij}\dot\lambda^i\dot\lambda^j}\,dt$~\cite{Crooks2007}. Within this slow-driving linear-response approximation, $W_{\mathrm{ex}}\ge\mathcal L^2/\tau$. Constant thermodynamic speed saturates the bound along a fixed path; minimizing also over paths gives a shortest geodesic. Where the normalized correlation integrals $\tau_{ij}$ are defined, $\zeta_{ij}=k_BT\,\tau_{ij}g_{ij}$ entry by entry. A common scalar relaxation time gives $\zeta=k_BT\tau_{\mathrm{relax}}g$ only when all the relevant entries have that same time~\cite{SivakCrooks2012}.


\section{Local tracking of a moving optimum}
\subsection{Tracking a moving optimum}
\label{subsec:lyapunov}

Let $\vartheta^*(t)$ be the minimizer of $L(\cdot,t)$, differentiable in $t$, and let $\|\dot\vartheta^*\|_g$ be its speed in the metric $g$. As a monitoring function take the squared geodesic distance
\begin{equation}
\label{eq:lyapunov}
V(t):=d_g\bigl(\vartheta(t),\vartheta^*(t)\bigr)^2 .
\end{equation}

\begin{theorem}[Loop tracking bound]
\label{thm:loop_stability}
Suppose the trajectory of~\eqref{eq:update_law} stays in a region where $d_g(\cdot,\vartheta^*(t))^2$ is smooth and $L(\cdot,t)$ is geodesically $m_L$-strongly convex for some $m_L>0$. Then
\begin{equation}
\label{eq:lyapunov_decrease}
\frac{dV}{dt}\ \le\ -2\eta m_L\,V+2\sqrt V\,\|\dot\vartheta^*\|_g .
\end{equation}
Consequently: (a) $V$ decreases whenever $\sqrt V>\|\dot\vartheta^*\|_g/(\eta m_L)$; (b) if $\|\dot\vartheta^*\|_g\le\Lambda$, then
\begin{equation}
\label{eq:tracking_neighbourhood}
\limsup_{t\to\infty}V(t)\ \le\ \frac{\Lambda^2}{(\eta m_L)^2};
\end{equation}
(c) if $\vartheta^*$ is constant, $V(t)\le V(0)\,e^{-2\eta m_Lt}$.
\end{theorem}

\begin{proof}
The natural gradient $g^{-1}\nabla L$ is the Riemannian gradient $\operatorname{grad}L$. At $\vartheta$ the gradient of $d_g(\cdot,\vartheta^*)^2$ is $-2\exp_\vartheta^{-1}\vartheta^*$. Strong convexity along the geodesic from $\vartheta$ to $\vartheta^*$ gives $L(\vartheta^*)\ge L(\vartheta)+\langle\operatorname{grad}L(\vartheta),\exp_\vartheta^{-1}\vartheta^*\rangle+\tfrac12m_LV$, and since $\vartheta^*$ is the minimizer, $L(\vartheta)\ge L(\vartheta^*)+\tfrac12m_LV$. Adding the two, $\langle\operatorname{grad}L(\vartheta),-\exp_\vartheta^{-1}\vartheta^*\rangle\ge m_LV$. The contribution of the update to $dV/dt$ is $\langle-2\exp_\vartheta^{-1}\vartheta^*,-\eta\operatorname{grad}L\rangle\le-2\eta m_LV$. The contribution of the moving optimum is $\langle-2\exp_{\vartheta^*}^{-1}\vartheta,\dot\vartheta^*\rangle\le2\sqrt V\|\dot\vartheta^*\|_g$ by Cauchy--Schwarz. This gives~\eqref{eq:lyapunov_decrease} and (a). For (b), $W=\sqrt V$ obeys $\dot W\le-\eta m_LW+\Lambda$ wherever $V>0$, so $\limsup W\le\Lambda/(\eta m_L)$. For (c) set $\dot\vartheta^*=0$.
\end{proof}

For the two-channel frame the loss $s^2+\rho^2\sin^2\theta$ has $L'(0)=0$ and $L''(0)=2\rho^2$, so its Riemannian Hessian at calibration is $2\rho^2/g(0)$. When probes make $g(0)>0$ (Proposition~\ref{prop:degenerate}), continuity gives strong convexity in a neighbourhood of calibration with $m_L$ close to $2\rho^2/g(0)$, and the theorem applies there. Without probes $g(0)=0$ and it does not apply.

Under drift the loop cannot converge to zero mismatch; it holds the mismatch within the neighbourhood~\eqref{eq:tracking_neighbourhood}, which shrinks with $\eta m_L$ and grows with the drift speed. Raising $\eta$ is not free: a physical implementation has to relate the adaptation rate to its energy budget (Section~\ref{sec:cost}), and in a stochastic loop a faster update also follows the noise, which is the trade-off Section~\ref{sec:price} prices exactly. A second constant governs how well the metric itself is conditioned.


\section{Thermodynamic cost}
\label{sec:cost}

\subsection{Cost components}
\label{subsec:cost_components}

A physical calibration loop senses, computes, actuates and prepares references. What each costs depends on how it is realized.

\emph{Sensing and computing.} Measurement, copying and logic can in principle be reversible, so reading a stream carries no positive cost by itself. If sensor and work registers irreversibly erase Shannon information at rates $R_s$ and $R_c$ bits per unit time, Landauer's bound gives $\dot W_{\mathrm{sense}}\ge k_BT\ln2\,R_s$ and $\dot W_{\mathrm{compute}}\ge k_BT\ln2\,R_c$.

\emph{References.} A reference read requires an input in a known state. Preparing a known pointer state from the thermal state of a degenerate pointer costs at least $k_BT\ln2$ per reference (Paper~I). At $R_{\mathrm{ref}}$ references per unit time this is $k_BT\ln2\,R_{\mathrm{ref}}$, unless references are recycled.

\emph{Actuation.} Moving the self-model from $\vartheta$ to a new estimate $\hat\vartheta^*$ is a finite-time transformation of whatever physical parameter stores it. In linear response, a protocol of duration $\tau_{\mathrm{recalib}}$ and thermodynamic length $\mathcal L(\vartheta,\hat\vartheta^*)$ dissipates at least $\mathcal L^2/\tau_{\mathrm{recalib}}$, a rate of at least $\mathcal L^2/\tau_{\mathrm{recalib}}^2$ if repeated every $\tau_{\mathrm{recalib}}$~\cite{SivakCrooks2012}.

\subsection{A conditional implementation bound}
\label{subsec:cost_theorem}

\begin{proposition}[Conditional implementation cost]
\label{prop:loop_cost_impl}
Suppose that (i) sensing, computing, reference preparation and actuation use separate physical degrees of freedom with no shared erasure accounting; (ii) sensor and work registers erase Shannon information at rates $R_s$ and $R_c$, and $R_{\mathrm{ref}}$ references per unit time are prepared from thermal states; and (iii) actuation follows a linear-response protocol of duration $\tau_{\mathrm{recalib}}$ and thermodynamic length $\mathcal L$, repeated every $\tau_{\mathrm{recalib}}$. Then
\begin{equation}
\label{eq:loop_cost}
\dot W_{\mathrm{loop}}\ \ge\ k_BT\ln2\,\bigl(R_s+R_c+R_{\mathrm{ref}}\bigr)+\frac{\mathcal L^2}{\tau_{\mathrm{recalib}}^2}.
\end{equation}
\end{proposition}

\begin{proof}
Landauer's bound applies to the two declared erasure streams and to each reference preparation; the Sivak--Crooks bound divided by the protocol duration gives the actuation rate. Separability allows the terms to be added without double counting. Dropping separability, or using reversible sensing and computation without resets in the window, removes the corresponding terms.
\end{proof}

The proposition and the drift law count different things. Remark~\ref{cor:thermo} counts references and the loss of acting on an imperfect self-model, in the exact Gaussian comparison. Proposition~\ref{prop:loop_cost_impl} adds the costs of the registers and actuators of one implementation. Together they bound what a physically realized calibration loop pays per use. Papers~I and~II supply no universal maintenance or processing cost that could be added to these without a concrete implementation, so no general persistence budget is claimed.

% ============================================================

\section{Finite-block margins and continuing filtering}
\subsection{A resource-conditioned operating margin}
\label{subsec:J_viability}

A lower bound on drive occupation cannot by itself supply an achievable gate. We instead choose the displaced squeezed Gaussian family of Q-RAIF~B and use its exact characteristic function for both error and energy. If its mean occupation is $N$, set
\[
V_Q=\frac1{2N+1},\quad
\lambda^2=\frac{\pi^2(2N+1)}{16N(N+1)},\quad
\alpha^2=\frac{\pi^2}{16\lambda^2},\quad
s=\frac{1+\exp(-2\lambda^2V_Q)}2.
\]
Then $\alpha^2+(V_Q+V_Q^{-1}-2)/4=N$. For the coherent comparison use $V_Q=1$, $\lambda^2=\pi^2/(16N)$ and the same formula for $s$. Each opportunity receives a fresh independent scheduled packet, costing $\theta_\omega N$ whether or not it is switched on. Actual switching adds $\theta_\omega\lambda^2$. Incoming field energy and switching work thus refer to the same state.

The gain is physical and bounded. In blocks of $L=4096$ food uses, draw independently one of $g_-=0.90h$ and $g_+=0.98h$, with equal prior probabilities, where $r=0.8^n$ and $h=\sqrt{1-r^2}$. Keep that gain constant through the block, with matched setting and known positive sign. At the start take $m=32$ references of the input $+1$. The count $K$ has likelihood
\[
p_g(k)=\binom mk[(1+g)/2]^k[(1-g)/2]^{m-k};
\qquad
\bar g_k=\frac{g_-p_{g_-}(k)+g_+p_{g_+}(k)}{p_{g_-}(k)+p_{g_+}(k)}.
\]
The reference pointer preparation costs $m\ln2$. We conservatively erase its outcomes separately for another $m\ln2$, and erase the $(m+1)$-state count register for $\ln(m+1)$ per block. This finite-memory filter incurs no Gaussian approximation or truncation.

Write
\[
S_r(x)=H_2\!\left(\frac{1+\sqrt{r^2+(1-r^2)x^2}}2\right).
\]
The posterior-optimal food levels give income $\ln2-H_2((1+s\bar g_k)/2)$. The food reset against the untouched quantum copy costs
$\ln2-S_r(0)+S_r(\bar g_k)$; the switching probability is exactly $1/2$. Its net, averaged over the exact law of $K$, is
\[
A=\E_K\!\left[S_r(0)-S_r(\bar g_K)-H_2((1+s\bar g_K)/2)\right]
-\theta_\omega(N+\lambda^2/2)-\frac{2m\ln2+\ln(m+1)}L.
\]
The reset is the reversible conditional-entropy limit for the stated classical--quantum state, retaining the copy marginal. Its attainability uses quasistatic processing across independent replicas of the complete block protocol, conditioned on the retained count. An implementation with reset overhead below a positive margin has the same sign of net work.

Economy~B adds one body read and drive opportunity per food use. The body starts each block uniform and is thermalized by a work-free swap with a fresh equilibrium bit at the block boundary. With bath flip probability $d=0.05$ and feedback on every negative read,
\[
x_{g,t+1}=0.9[(1-s)x_{g,t}+sg],\qquad x_{g,0}=0.
\]
Given $K=k$, define $x_k=\E[x_{g,t}\mid k]$, $y_k=\E[gx_{g,t}\mid k]$, and $p_\pm=(1\pm y_k)/2$. The conditional body-copy biases are $b_\pm=(x_k\pm\bar g_k)/(1\pm y_k)$. Body erasure therefore costs
\[
H_2(p_+)-S_r(x_k)+\sum_{\pm}p_\pm S_r(b_\pm),
\]
and its drive costs $\theta_\omega[N+p_-\lambda^2]$. Subtract their exact block average from $A$ to get $B$. This policy retains only $K$, ignores body readings as calibration data, and requires no unproved body-stream information rate.

\begin{table}[t]
\centering\small
\begin{tabular}{@{}rrlrrrrr@{}}
\toprule
$n$ & $\theta_\omega$ & Drive & $N$ & A: margin & B: margin & Body loss & B: purity\\
\midrule
3 & 0.01 & coherent & 6.220 & -0.1607 & -0.5798 & 0.1402 & 8.334 \\
3 & 0.01 & squeezed & 4.220 & -0.0910 & -0.4899 & 0.1377 & 8.334 \\
5 & 0.01 & coherent & 7.343 & 0.0533 & -0.2522 & 0.1037 & 13.890 \\
5 & 0.01 & squeezed & 4.758 & 0.1393 & -0.1402 & 0.1013 & 13.890 \\
7 & 0.01 & coherent & 7.865 & 0.1552 & -0.0985 & 0.0896 & 19.446 \\
7 & 0.01 & squeezed & 5.019 & 0.2492 & 0.0240 & 0.0873 & 19.446 \\
5 & 0.05 & coherent & 2.522 & -0.1104 & -0.4709 & 0.1097 & 13.890 \\
5 & 0.05 & squeezed & 2.454 & 0.0036 & -0.3540 & 0.1032 & 13.890 \\
7 & 0.05 & coherent & 2.708 & -0.0196 & -0.3317 & 0.0955 & 19.446 \\
7 & 0.05 & squeezed & 2.587 & 0.1069 & -0.1999 & 0.0891 & 19.446 \\
10 & 0.10 & coherent & 1.587 & -0.0810 & -0.3884 & 0.0945 & 27.780 \\
10 & 0.10 & squeezed & 1.919 & 0.0410 & -0.3007 & 0.0840 & 27.780 \\
\bottomrule
\end{tabular}
\caption{Achievable conditional-entropy limit in units of $k_BT$. Margins are harvest minus the listed reset, reference and drive costs after food and record resources have been supplied. Each row optimizes food-only drive within its family, then uses that same state for food and body. The B: purity column is external record-memory purity per food use, $(4n+mn/L)\ln2$, in addition to the listed operating costs; A uses $(2n+mn/L)\ln2$. Body loss is averaged after each step over the full block. Calibration costs $0.011684$ per food use.}
\label{tab:viability}
\end{table}

At $n=7$ and $\theta_\omega=0.01$, the squeezed construction leaves $0.02396\,k_BT$ per food use after body upkeep, with body loss $0.08725$; the coherent comparison leaves $-0.09847\,k_BT$. A sufficiently small irreversible overhead preserves the positive squeezed margin. This is a maintained body supplied by a record and drive stream, with a gain that changes between blocks. It is not an infinite Gaussian-gain process.

Record formation consumes $2n\ln2$ of fresh memory purity per food use in~A and $4n\ln2$ in~B, plus $mn\ln2/L$ for reference fragments. Those memories are supplied by the environment and their purity is explicitly an input, not work produced by the cycle. Replenishing them from the agent's own thermal resources debits that input and eliminates the apparent profit. The drive additionally presupposes its preparation, phase reference and scheduled port; uncomputed apparatus overhead can be tolerated only within the displayed positive margin. The example identifies a nonempty region of positive operating margin under these resource declarations. It does not derive the resources or the reference standard from the loop.

\subsection{A continuing joint filter and adaptive references}
\label{subsec:online_joint}

The block calculation redraws a prior and resets the body at known times. A continuing observer instead carries its model through the changes that its own history reveals. Keep the $n=7$ fragments and the same single-pass squeezed drive, $N=5.01898$, $s=0.979992$ and $\lambda^2=0.225388$. Let the gain switch between $g_-=0.90h$ and $g_+=0.98h$ with probability $\delta=3\times10^{-4}$ per use, and the sign invert with probability $\eta=2\times10^{-4}$. The body bath has $d=0.05$. These are bounded transitions valid for arbitrarily long operation. The body, instrument and stored model continue through every reference-budget boundary.

The hidden state $Z_t=(G_t,\sigma_t,B_t)$ has eight values. Each use gives one body reading with likelihood $(1+y\sigma g B)/2$. A requested reference gives one known-input reading with likelihood $(1+z\sigma g)/2$. Its request can depend on the body reading and earlier records. The observer then selects a body flip by the sign of its posterior body mean, and food levels by its posterior signed gain. It receives two fresh scheduled drive packets per use, one for each task, and a maximum of $m=32$ reference ports per $L=4096$ uses. The ports, their timing and their input labels are supplied equally to the compared policies; an unused port creates no additional income.

For body command $a\in\{0,1\}$, let $T_a$ be the eight-state transition matrix. Its gain and sign factors are the symmetric switching kernels of rates $\delta$ and $\eta$. The body flips with probability $d$ for $a=0$, and with probability $d+s(1-2d)$ for $a=1$. All entries of $T_a$ are positive. Given the selected actions and realized records, the filter predicts with $T_a$ and multiplies by the two likelihoods, omitting the second when no reference is requested. The request rule contributes no extra parameter-dependent likelihood factor. The fair food reading alone supplies no instrument information and is used only for its task.

Positive transition and observation kernels give a uniform finite-register
certificate for this filter (Proposition~\ref{prop:finite_filter},
Appendix~\ref{app:filter_certificates}). It controls quantization error
even when actions and reference requests are chosen from the rounded
posterior.

This filter forms its next prior from the preceding posterior and the actual transition law. It does not request a newly assigned gain prior or reprepare the body. For the stated parameters, $a_0=3\times10^{-9}$, $D=39.14572$ and $\tau=0.999999993681$. A log grid with $\Delta=2^{-44}$ uses 350 model bits; its asymptotic total-variation bound is $2.25\times10^{-6}$ and its one-step excess-risk bound is $4.50\times10^{-6}$. These conservative uniform bounds concern quantization of a correct filter. They do not certify a wrong transition family or the arithmetic implementation. The 350-bit state is an abstract finite controller; its workspace, update throughput and erasure protocol remain to be implemented.

For a resource-matched comparison, the periodic policy uses one reference every 128 steps. The adaptive policy requests a reference whenever the posterior probability of the less likely sign exceeds $0.04$, subject to the same 32-port cap. Both update the joint model from every body reading and use the same body-control rule. Starting with the uniform eight-state prior, 128 independent trajectories of 65536 steps, with the first 8192 discarded, give
\[
\begin{array}{lccc}
 &\text{Body loss}&\text{References/use}&\text{Food mismatch}\\
\text{Periodic}&0.11726\pm0.00083&0.0078125&0.06726\pm0.00048\\
\text{Adaptive}&0.10319\pm0.00066&0.005704\pm0.000020&0.05226\pm0.00014 .
\end{array}
\]
Errors are one standard error across independent trajectories. Food mismatch is the posterior Jensen gap $H_2((1+s|\E[\sigma g\mid\mathcal H_t]|)/2)-\E[H_2((1+sg)/2)\mid\mathcal H_t]$ in nats. The adaptive gain mean has mean-square error $0.000353$, compared with the stationary prior variance $0.001529$. A separate seed and shorter horizon give body losses $0.11730$ and $0.10344$. Requesting too early is costly under the cap: threshold $0.01$ exhausts references prematurely and gives body loss $0.19567$ in that shorter run. Adaptive timing needs an allocation rule as well as a confidence statistic.

Every policy is allocated the same food and record stream, drive state and 32-port reference budget. The external record purity is at most $(4n+mn/L)\ln2=19.44603$ nats per food use, and reference preparation plus conservative separate erasure is charged at the same supplied-cap rate $2m\ln2/L$. The posterior is updated from the body's own records and the policy acts on that body, closing the inference--control loop. This establishes continuing model revision and an operational benefit at finite memory. It does not close the apparatus resource loop: conditional-copy reset, model-register processing, field preparation, phase reference and the clock still need one attainable cyclic implementation. The block margin of Table~\ref{tab:viability} cannot be imported into this continuing protocol without that ledger.


\section{Finite-filter certificates}\label{app:filter_certificates}

\subsection{Continuing eight-state filter}

\begin{proposition}[Finite continuing filter]
\label{prop:finite_filter}
Let $a_0=\min_{a,i,j}(T_a)_{ij}>0$, and let every observation likelihood lie in $[l_-,l_+]$ with $l_->0$. Suppose at most two observations precede each filter update. Set
\[
 \epsilon_0=a_0(l_-/l_+)^2,\quad A=-\ln\epsilon_0,\quad
 D=\max_{a,i,j,k,l}\ln\frac{(T_a)_{ik}(T_a)_{jl}}{(T_a)_{il}(T_a)_{jk}},
 \qquad \tau=\tanh(D/4)<1.
\]
After the first prediction, every exact update has probability at least $\epsilon_0$ in each state. Round each of the seven log ratios $\ln(p_i/p_8)$ to a grid of spacing $\Delta$, then normalize their exponentials. This uses at most
\[
 7\left\lceil\log_2\left(\frac{2A}{\Delta}+3\right)\right\rceil
\]
bits for the model. For the exact posterior $p_t^*$ conditioned on the same realized history and the rounded filter $\tilde p_t$,
\[
 d_H(\tilde p_t,p_t^*)\le\tau^t d_H(\tilde p_0,p_0^*)+
       \frac{\Delta(1-\tau^t)}{1-\tau},\qquad
 \|\tilde p_t-p_t^*\|_{\rm TV}\le\tanh\bigl(d_H(\tilde p_t,p_t^*)/4\bigr).
\]
Here $d_H(p,q)=\max_i\ln(p_i/q_i)-\min_i\ln(p_i/q_i)$. A one-step action minimizing a loss in $[0,1]$ under $\tilde p_t$ has excess conditional risk at most $2\|\tilde p_t-p_t^*\|_{\rm TV}$. With a finite reference counter and clock, the eight-state plant and rounded controller form a finite Markov chain and possess an invariant distribution under the supplied streams.
\end{proposition}

\begin{proof}
Prediction gives each probability at least $a_0$. The product of at most two positive likelihoods reduces that bound by at most $(l_-/l_+)^2$. Hence the seven log ratios lie in $[-A,A]$, and rounding increases their range by at most $\Delta/2$ at each end. In the Hilbert metric, positive diagonal likelihood multiplication and normalization are isometries, and multiplication by $T_a$ contracts by at most $\tau$, the Birkhoff coefficient. This is the standard projective stability mechanism for finite hidden Markov filters~\cite{LeGlandOudjane2004}. Rounding changes the oscillation of the log ratios by at most $\Delta$. The triangle inequality gives $e_t\le\tau e_{t-1}+\Delta$, which iterates to the displayed bound even when actions and reference requests depend on earlier rounded states. To obtain the total-variation constant, let $m\le p_i/q_i\le M$, with $M/m=e^e$. Normalization and the chord bound for $(x-1)_+$ give
\[
 \|p-q\|_{\rm TV}\le\frac{(M-1)(1-m)}{M-m}\le\tanh(e/4),
\]
where maximizing over $m$ at fixed $M/m$ gives $m=e^{-e/2}$; distance zero is immediate. Comparing a minimizing action under each distribution costs at most twice that distance for a bounded loss. The model grid, budget counter, clock and eight-state plant are finite, so their joint transition has an invariant probability. This last statement need not give a unique invariant distribution or an optimal reference policy, and it concerns the declared plant and controller rather than the whole resource-supplying apparatus.
\end{proof}

\subsection{Filter on working invocations}

The hard access gate has zero likelihoods on excluded body states, so Proposition~\ref{prop:finite_filter}'s eight-state certificate cannot simply be inherited. Filtering at actual working invocations avoids this issue. Each of the two event matrices is positive. For the present switching rates,
\[
 \min_{\ell,i,j}(\mathsf K_\ell)_{ij}\ge
 a=(1-d)\delta\eta=\frac{57}{10^9},\qquad g_{\max}<0.98.
\]
One body likelihood multiplication gives a posterior lower bound $a/99$, hence a log-ratio range less than $22$. Event cross-ratios are at most $a^{-2}$, so their common projective contraction coefficient is at most $(1-a)/(1+a)$. Round the three log ratios on a grid $\Delta=2^{-44}$. The proof of Proposition~\ref{prop:finite_filter}, now applied once per working invocation, gives 150 ideal model bits and
\[
 \limsup\|\tilde p_n-p_n^*\|_{\rm TV}
 \le\frac{\Delta(1+a)}{8a}<1.247\times10^{-7}.
\]
For a one-step loss in $[0,1]$, regret is at most twice the total-variation distance and hence less than $2.494\times10^{-7}$. The latch, startup state, arithmetic, workspace and protection are separate from these model bits. In wall time the plant, held finite register and latch form a finite Markov chain; they have an invariant probability under the supplied streams. The claim neither processes new data during disability nor certifies a cyclic processor.

\begin{thebibliography}{99}

\bibitem{Liu2026TDOME_I}
S.~Liu, \emph{Memory as a Usable Resource: Correlations, Records and Complete Endpoint Accounting} (T-DOME~I), 2026; July 2026 record DOI: 10.5281/zenodo.18574342.

\bibitem{Liu2026TDOME_II}
S.~Liu, \emph{Reference Frames Selected by History: The Capability and Cost of a Bounded Agent's Self-Frame} (T-DOME~II), 2026; July 2026 record DOI: 10.5281/zenodo.18579703.

\bibitem{Liu2026CIII}
S.~Liu, \emph{Finite Self-Models and the Explanatory Gap: What a Continuing Observer Can Know About Itself} (Consciousness~III), 2026; July 2026 record DOI: 10.5281/zenodo.19607724.

\bibitem{Liu2026HAFF_C}
S.~Liu, \emph{Causation, Agency, and Existence: Historical Control and Maintained Intervention}, Zenodo (2026), DOI: 10.5281/zenodo.18374805.

\bibitem{Liu2026HAFF_E}
S.~Liu, \emph{Measurement as Accessibility: Measurement Permissions, Stable Records and Public Readout}, revised version, Zenodo (2026), DOI: 10.5281/zenodo.18400065.

\bibitem{Liu2026HAFF_F}
S.~Liu, \emph{Records, Dissipation and Intervention: Conditions for a Temporal Direction}, revised version, Zenodo (2026), DOI: 10.5281/zenodo.18400425.

\bibitem{Liu2026Chain}
S.~Liu, \emph{The Imprint Chain: Records, Self-Calibration, and Physical Feedback}, (2026), series record DOI: 10.5281/zenodo.20151594.

\bibitem{Liu2026HAFF_G}
S.~Liu, \emph{Internal Selection and Revisable Organization: How History Changes Permissions and Physical Sources} (HAFF Paper~G), 2026; July 2026 record DOI: 10.5281/zenodo.18402907.

\bibitem{Liu2026QRAIF_B}
S.~Liu, \emph{Thermodynamic Stability Constraints on Persistent Open Quantum Subsystems: Control Algebra, Physical Actuation and Energy Cost} (Q-RAIF~B), 2026; series record DOI: 10.5281/zenodo.18525890.

\bibitem{Liu2026FII}
S.~Liu, \emph{Foundation II: Shared Manifestation. Shared Structure, Public Calibration and Different Histories}, 2026; series record DOI: 10.5281/zenodo.19641372.

\bibitem{Liu2026FIII}
S.~Liu, \emph{Foundation III: The Perfuming Loop. Held Content, Learned Frames and Physical Return}, 2026; July 2026 record DOI: 10.5281/zenodo.19642189.

\bibitem{Liu2026BRAIN}
S.~Liu, \emph{The Brain as a Plastic Observable-Algebra Network:
Learning Shared Access from Sensorimotor History} (Brain Anchor),
2026; series concept DOI: 10.5281/zenodo.20044329.

\bibitem{Rao1945}
C.~R.~Rao,
\emph{Information and the accuracy attainable in the estimation of statistical parameters},
Bull.\ Calcutta Math.\ Soc.\ \textbf{37}, 81 (1945).

\bibitem{Cencov1982}
N.~N.~\v{C}encov,
\emph{Statistical Decision Rules and Optimal Inference},
Translations of Mathematical Monographs \textbf{53}, AMS (1982).

\bibitem{Amari1998}
S.-i.~Amari,
\emph{Natural gradient works efficiently in learning},
Neural Computation \textbf{10}, 251 (1998).

\bibitem{AmariNagaoka2000}
S.-i.~Amari and H.~Nagaoka,
\emph{Methods of Information Geometry},
Translations of Mathematical Monographs \textbf{191}, AMS (2000).

\bibitem{vanTrees1968}
H.~L.~van Trees,
\emph{Detection, Estimation, and Modulation Theory}, Part~I, Wiley (1968).

\bibitem{Tichavsky1998}
P.~Tichavsk\'y, C.~H.~Muravchik and A.~Nehorai,
\emph{Posterior Cram\'er--Rao bounds for discrete-time nonlinear filtering},
IEEE Trans.\ Signal Process.\ \textbf{46}, 1386--1396 (1998).

\bibitem{Kalman1960}
R.~E.~Kalman,
\emph{A new approach to linear filtering and prediction problems},
J.\ Basic Eng.\ \textbf{82}, 35--45 (1960).

\bibitem{LeGlandOudjane2004}
F.~Le Gland and N.~Oudjane,
\emph{Stability and uniform approximation of nonlinear filters using the Hilbert metric, and application to particle filters},
Ann.\ Appl.\ Probab.\ \textbf{14}, 144--187 (2004), DOI: 10.1214/aoap/1075828050.

\bibitem{AstromWittenmark1995}
K.~J.~\AA{}str\"{o}m and B.~Wittenmark,
\emph{Adaptive Control}, 2nd ed., Addison-Wesley (1995).

\bibitem{Crooks2007}
G.~E.~Crooks,
\emph{Measuring thermodynamic length},
Phys.\ Rev.\ Lett.\ \textbf{99}, 100602 (2007).

\bibitem{SivakCrooks2012}
D.~A.~Sivak and G.~E.~Crooks,
\emph{Thermodynamic metrics and optimal paths},
Phys.\ Rev.\ Lett.\ \textbf{108}, 190602 (2012).

\bibitem{BaratoSeifert2015}
A.~C.~Barato and U.~Seifert,
\emph{Thermodynamic uncertainty relation for biomolecular processes},
Phys.\ Rev.\ Lett.\ \textbf{114}, 158101 (2015).

\bibitem{Ito2018}
S.~Ito,
\emph{Stochastic thermodynamic interpretation of information geometry},
Phys.\ Rev.\ Lett.\ \textbf{121}, 030605 (2018).

\bibitem{Ashby1956}
W.~R.~Ashby,
\emph{An Introduction to Cybernetics}, Chapman \& Hall (1956).

\bibitem{vonFoerster2003}
H.~von Foerster,
\emph{Understanding Understanding: Essays on Cybernetics and Cognition}, Springer (2003).

\bibitem{CoverThomas2006}
T.~M.~Cover and J.~A.~Thomas,
\emph{Elements of Information Theory}, 2nd ed., Wiley (2006).

\end{thebibliography}

\end{document}
